% Options for packages loaded elsewhere
\PassOptionsToPackage{unicode}{hyperref}
\PassOptionsToPackage{hyphens}{url}
%
\documentclass[
]{article}
\usepackage{amsmath,amssymb}
\usepackage{iftex}
\ifPDFTeX
  \usepackage[T1]{fontenc}
  \usepackage[utf8]{inputenc}
  \usepackage{textcomp} % provide euro and other symbols
\else % if luatex or xetex
  \usepackage{unicode-math} % this also loads fontspec
  \defaultfontfeatures{Scale=MatchLowercase}
  \defaultfontfeatures[\rmfamily]{Ligatures=TeX,Scale=1}
\fi
\usepackage{lmodern}
\ifPDFTeX\else
  % xetex/luatex font selection
\fi
% Use upquote if available, for straight quotes in verbatim environments
\IfFileExists{upquote.sty}{\usepackage{upquote}}{}
\IfFileExists{microtype.sty}{% use microtype if available
  \usepackage[]{microtype}
  \UseMicrotypeSet[protrusion]{basicmath} % disable protrusion for tt fonts
}{}
\makeatletter
\@ifundefined{KOMAClassName}{% if non-KOMA class
  \IfFileExists{parskip.sty}{%
    \usepackage{parskip}
  }{% else
    \setlength{\parindent}{0pt}
    \setlength{\parskip}{6pt plus 2pt minus 1pt}}
}{% if KOMA class
  \KOMAoptions{parskip=half}}
\makeatother
\usepackage{xcolor}
\usepackage{longtable,booktabs,array}
\usepackage{calc} % for calculating minipage widths
% Correct order of tables after \paragraph or \subparagraph
\usepackage{etoolbox}
\makeatletter
\patchcmd\longtable{\par}{\if@noskipsec\mbox{}\fi\par}{}{}
\makeatother
% Allow footnotes in longtable head/foot
\IfFileExists{footnotehyper.sty}{\usepackage{footnotehyper}}{\usepackage{footnote}}
\makesavenoteenv{longtable}
\setlength{\emergencystretch}{3em} % prevent overfull lines
\providecommand{\tightlist}{%
  \setlength{\itemsep}{0pt}\setlength{\parskip}{0pt}}
\setcounter{secnumdepth}{-\maxdimen} % remove section numbering
\ifLuaTeX
  \usepackage{selnolig}  % disable illegal ligatures
\fi
\IfFileExists{bookmark.sty}{\usepackage{bookmark}}{\usepackage{hyperref}}
\IfFileExists{xurl.sty}{\usepackage{xurl}}{} % add URL line breaks if available
\urlstyle{same}
\hypersetup{
  hidelinks,
  pdfcreator={LaTeX via pandoc}}

\author{}
\date{}

\begin{document}

\hypertarget{nasco-feeding-minds-estrategia-inbound-basada-en-el-diagnuxf3stico}{%
\section{NASCO Feeding Minds: estrategia Inbound basada en el
diagnóstico}\label{nasco-feeding-minds-estrategia-inbound-basada-en-el-diagnuxf3stico}}

\hypertarget{ea1-entrega-final}{%
\subsection{EA1: Entrega final}\label{ea1-entrega-final}}

\textbf{Escenario de Aprendizaje 1}\\
\textbf{Asignatura:} Inteligencia Artificial y Marketing Digital\\
\textbf{Curso:} 2026-2027\\
\textbf{Autor:} Sergio Morales Cobo

\begin{center}\rule{0.5\linewidth}{0.5pt}\end{center}

\hypertarget{resumen-ejecutivo}{%
\subsection{Resumen ejecutivo}\label{resumen-ejecutivo}}

NASCO Feeding Minds es una ONG que une educación digital y oportunidades
de desarrollo en Ghana. Su propuesta va dirigida a particulares (B2C) y
empresas (B2B) que quieren contribuir a una causa educativa, ya sea con
dinero o con material informático. La memoria de 2024 dice que necesitan
estabilizar la financiación, actualizar la web, captar socios y reforzar
las redes sociales. Así que el reto de Inbound no es solo ganar
notoriedad, sino transformar el interés por la causa en confianza y
colaboración que se mantenga en el tiempo.
\href{https://nascoict.org/wp-content/uploads/MEMORIA-2024.pdf}{F03}

Mirando sus redes sociales he visto que tienen unos 34.000 seguidores en
Instagram, 5.790 en Facebook, 2.415 en LinkedIn y 2.400 en X. Tienen
proyectos documentados y las cuatro cuentas están enlazadas desde la web
oficial. \href{https://www.instagram.com/nasco_ict/}{F08}
\href{https://www.facebook.com/NascoICT/}{F09}
\href{https://twitter.com/nasco_ict/}{F11}
\href{https://www.linkedin.com/posts/nasco-feeding-minds_nascofoundation-educaci\%C3\%B3ndigital-impactosocial-activity-7500847528586846208-5q6E}{F12}

La debilidad más visible es que la información no es consistente: la
portada dice 11 aulas, la página de aulas 17 y ``Quiénes somos'' 22, sin
explicar a qué periodo corresponde cada cifra. Es posible que 22 sea el
total acumulado desde 2012 y que 11 o 17 sean cifras de un momento dado,
pero ninguna página lo aclara. Tampoco cuadran los alumnos: la portada
habla de 66.000 personas y 3.766 alumnos actuales, ``Quiénes somos'' de
60.000 estudiantes, y la página de aulas de 6.600 y 6.000 dentro de la
misma página. No creo que los datos sean falsos, pero así presentados es
difícil hacerse una idea clara y actualizada de la marca.
\href{https://nascoict.org/}{F01}
\href{https://nascoict.org/quienes-somos/}{F02}
\href{https://nascoict.org/aulas-informaticas/}{F04}

Mi primera prioridad es que más personas en España descubran NASCO,
tanto posibles socios particulares como empresas colaboradoras. Al
revisar la página de colaboración me ha parecido claro cómo donar o
hacerse socio, así que no lo planteo como el problema principal. Las
cifras de seguidores muestran que NASCO ya tiene presencia digital, pero
no permiten saber a cuántas personas nuevas llega ni cuántas colaboran:
para eso harían falta datos internos. La coherencia de la información y
los resultados educativos siguen siendo importantes para sostener la
confianza de quien llegue a conocer la organización.

\begin{center}\rule{0.5\linewidth}{0.5pt}\end{center}

\hypertarget{alcance-muxe9todo-y-criterios-de-evidencia}{%
\subsection{1. Alcance, método y criterios de
evidencia}\label{alcance-muxe9todo-y-criterios-de-evidencia}}

\hypertarget{alcance-de-esta-entrega}{%
\subsubsection{1.1. Alcance de esta
entrega}\label{alcance-de-esta-entrega}}

Este archivo empezó como la entrega intermedia de Marca y Entorno y lo
he ampliado hasta el EA1 final, añadiendo recursos, competencia, dos
buyer persona, customer journey, brújula estratégica, objetivos y
medición. La prioridad que he elegido tras revisar el borrador es
\textbf{ampliar el conocimiento de NASCO entre particulares y empresas
en España}. No parto de que la página de colaboración sea confusa,
porque al revisarla me parece clara. Esto es un plan académico basado en
fuentes públicas: no representa un plan aprobado por NASCO, pero las
decisiones de análisis y priorización son mías.

\textbf{Mercado de captación:} España.\\
\textbf{Territorio de intervención:} Ghana.\\
\textbf{Horizonte propuesto para los objetivos:} seis meses desde la
puesta en marcha, condicionado a que haya responsables y acceso a
medición.

\hypertarget{muxe9todo}{%
\subsubsection{1.2. Método}\label{muxe9todo}}

He utilizado:

\begin{enumerate}
\def\labelenumi{\arabic{enumi}.}
\tightlist
\item
  Material docente de la asignatura.
\item
  Web y memoria de actividad de NASCO.
\item
  Perfiles sociales enlazados desde la web oficial y una muestra de
  publicaciones.
\item
  Fuentes institucionales: INE, Banco Mundial, Comisión Europea, BOE,
  ITU/UNITAR, AEPD y AEAT.
\item
  Fuentes sectoriales: Asociación Española de Fundraising y contenidos
  públicos de Consumer Trends 2026 de Zorraquino.
\end{enumerate}

Las citas \textbf{F01-F14} corresponden a la web y redes sociales de
NASCO; \textbf{F15-F28} a fuentes institucionales y sectoriales del
entorno (PESTEL y tendencias); \textbf{F29-F30} a páginas legales y
versión inglesa de la web; \textbf{F31-F37} a las alternativas y
alianzas consultadas. He comprobado personalmente la totalidad de las
fuentes del dossier: las cifras de aulas, los seguidores en redes, la
página de colaboración y la memoria 2024.

\hypertarget{cuxf3mo-interpreto-los-resultados}{%
\subsubsection{1.3. Cómo interpreto los
resultados}\label{cuxf3mo-interpreto-los-resultados}}

A lo largo del dossier distingo entre:

\begin{itemize}
\tightlist
\item
  \textbf{Lo que vi:} datos que recuperé de una página o perfil
  identificado.
\item
  \textbf{Lo que NASCO dice de sí misma:} información institucional que
  no equivale a una evaluación independiente del impacto.
\item
  \textbf{Lo que yo interpreto:} significado estratégico que razono a
  partir de la evidencia.
\item
  \textbf{Lo que propongo:} recomendaciones que todavía no están
  validadas con datos internos o entrevistas.
\end{itemize}

\begin{center}\rule{0.5\linewidth}{0.5pt}\end{center}

\hypertarget{presentaciuxf3n-de-la-marca}{%
\subsection{2. Presentación de la
marca}\label{presentaciuxf3n-de-la-marca}}

\hypertarget{actividad-e-identidad}{%
\subsubsection{2.1. Actividad e identidad}\label{actividad-e-identidad}}

NASCO Feeding Minds pertenece al tercer sector y a la cooperación al
desarrollo, especializada en educación y alfabetización digital. La
memoria de 2024 la describe como una asociación fundada en 2012 por
Ousman Umar y Banasco Seidu. Su actividad consiste en aulas informáticas
y programas de formación en Ghana.
\href{https://nascoict.org/wp-content/uploads/MEMORIA-2024.pdf}{F03}
(pp.~3-5)

La misma memoria dice que durante 2024 ofreció educación digital a
\textbf{58 escuelas y más de 6.000 alumnos a través de 17 aulas}. Sin
embargo, la web no es consistente: la portada dice 11 aulas, ``Quiénes
somos'' dice 22, y los alumnos también varían según la página (66.000,
60.000, 6.600, 6.000 y 3.766). Es posible que 22 sea el total acumulado
desde 2012, pero ninguna página lo aclara.
\href{https://nascoict.org/}{F01}
\href{https://nascoict.org/quienes-somos/}{F02}
\href{https://nascoict.org/wp-content/uploads/MEMORIA-2024.pdf}{F03}
\href{https://nascoict.org/aulas-informaticas/}{F04}

La web también utiliza el nombre ``NASCO Foundation'' en la portada,
aunque el logo y el pie de página dicen ``NASCO Feeding Minds''. En
LinkedIn también aparece ``NASCO Foundation'' en algunas publicaciones.
He buscado en la web, en la memoria 2018 y en perfiles públicos y no he
encontrado ningún documento que explique si son la misma entidad con dos
nombres, una traducción o un ``rebranding''. En este trabajo mantengo
``NASCO Feeding Minds'' como nombre de referencia.
\href{https://nascoict.org/}{F01}
\href{https://nascoict.org/quienes-somos/}{F02}
\href{https://www.linkedin.com/posts/nasco-feeding-minds_nascofoundation-educaci\%C3\%B3ndigital-impactosocial-activity-7500847528586846208-5q6E}{F12}

\hypertarget{puxfablicos-y-relaciuxf3n-b2bb2c}{%
\subsubsection{2.2. Públicos y relación
B2B/B2C}\label{puxfablicos-y-relaciuxf3n-b2bb2c}}

\begin{longtable}[]{@{}
  >{\raggedright\arraybackslash}p{(\columnwidth - 4\tabcolsep) * \real{0.3333}}
  >{\raggedright\arraybackslash}p{(\columnwidth - 4\tabcolsep) * \real{0.3333}}
  >{\raggedright\arraybackslash}p{(\columnwidth - 4\tabcolsep) * \real{0.3333}}@{}}
\toprule\noalign{}
\begin{minipage}[b]{\linewidth}\raggedright
Público
\end{minipage} & \begin{minipage}[b]{\linewidth}\raggedright
Relación con NASCO
\end{minipage} & \begin{minipage}[b]{\linewidth}\raggedright
Papel en el proyecto
\end{minipage} \\
\midrule\noalign{}
\endhead
\bottomrule\noalign{}
\endlastfoot
Estudiantes y comunidades de Ghana & Beneficiarios de los programas. &
Personas en quienes se busca generar valor educativo y social. \\
Particulares residentes en España & Donantes puntuales y socios. &
Público de captación B2C, por analogía con una relación
organización-particular. \\
Empresas y entidades en España & Colaboradoras económicas o donantes de
equipos. & Público de captación B2B, por analogía con una relación entre
organizaciones. \\
\end{longtable}

NASCO no se presenta como una pyme mercantil, ya que, es una ONG. Los
términos B2B y B2C los adapto al fundraising: \textbf{colaborar no
equivale a comprar, y el donante no se debe confundir con el
beneficiario}. Las vías para particulares y empresas se ven en las
páginas de colaboración y donación de equipos.
\href{https://nascoict.org/aulas-informaticas/}{F04}
\href{https://nascoict.org/hazte-socio/}{F05}

\hypertarget{motivo-de-elecciuxf3n-y-reto-de-marketing}{%
\subsubsection{2.3. Motivo de elección y reto de
marketing}\label{motivo-de-elecciuxf3n-y-reto-de-marketing}}

He elegido NASCO porque reúne tres ámbitos que me interesan: la
informática, la educación y la psicología. El proyecto me permite
analizar cómo explicar una iniciativa educativa y qué necesita saber una
persona antes de decidir si colabora. Además, al ser una organización
real, hay proyectos, canales y documentos públicos sobre los que
trabajar.

La memoria de 2024 plantea actualizar la web, ejecutar campañas de
captación y desarrollar la presencia en redes. También recoge una
aspiración de \textbf{2.000 socios}.
\href{https://nascoict.org/wp-content/uploads/MEMORIA-2024.pdf}{F03}
(pp.~13-15)

\begin{center}\rule{0.5\linewidth}{0.5pt}\end{center}

\hypertarget{cuxedrculo-de-oro-y-relato-de-marca}{%
\subsection{3. Círculo de Oro y relato de
marca}\label{cuxedrculo-de-oro-y-relato-de-marca}}

\hypertarget{aplicaciuxf3n-del-modelo-de-sinek}{%
\subsubsection{3.1. Aplicación del modelo de
Sinek}\label{aplicaciuxf3n-del-modelo-de-sinek}}

Aplico el Círculo de Oro de Sinek a lo que he visto en la web y la
memoria de NASCO:

\begin{longtable}[]{@{}
  >{\raggedright\arraybackslash}p{(\columnwidth - 2\tabcolsep) * \real{0.5000}}
  >{\raggedright\arraybackslash}p{(\columnwidth - 2\tabcolsep) * \real{0.5000}}@{}}
\toprule\noalign{}
\begin{minipage}[b]{\linewidth}\raggedright
Nivel
\end{minipage} & \begin{minipage}[b]{\linewidth}\raggedright
Qué veo en NASCO
\end{minipage} \\
\midrule\noalign{}
\endhead
\bottomrule\noalign{}
\endlastfoot
\textbf{PORQUÉ} (el propósito) & El lugar de nacimiento no debería
impedir a una persona construir un futuro digno. Los jóvenes tienen que
tener oportunidades para decidir su trayectoria y migrar debería poder
ser una elección, no una salida forzada por la falta de alternativas. \\
\textbf{CÓMO} (lo que hacen) & 1. Acercan formación digital a
comunidades y centros educativos de Ghana. 2. Combinan equipamiento con
formación, en lugar de limitarse a entregar dispositivos. 3. Aprovechan
equipos informáticos aportados por colaboradores. 4. Cooperan con
personas, empresas y entidades para desarrollar proyectos concretos. 5.
Publican memorias e información de actividad para explicar lo
realizado. \\
\textbf{QUÉ} (lo que ofrecen) & Aulas de informática; formación digital;
programas de formación de formadores; vías de donación y colaboración
periódica; canal de aportación de material informático. \\
\end{longtable}

\textbf{Evidencias del CÓMO:} memoria de actividad; página de aulas, que
describe formación y reutilización de equipos; sección de transparencia;
publicaciones de proyectos con Barça Foundation y Deloitte.
\href{https://nascoict.org/wp-content/uploads/MEMORIA-2024.pdf}{F03}
\href{https://nascoict.org/aulas-informaticas/}{F04}
\href{https://nascoict.org/transparencia/}{F06}
\href{https://www.linkedin.com/posts/nasco-feeding-minds_nascofoundation-feedingminds-digitaleducation-activity-7465326634426843136-PKBW}{F13}
\href{https://www.linkedin.com/posts/nasco-feeding-minds_nasco-x-deloitte-activity-7464593630444306432-8i9r}{F14}

\hypertarget{historia-frente-a-relato}{%
\subsubsection{3.2. Historia frente a
relato}\label{historia-frente-a-relato}}

La historia del fundador explica el origen de la organización, pero el
relato estratégico tiene que ir más allá de esa biografía: \textbf{crear
capacidades y oportunidades mediante educación digital}. Esto permite
que la marca conserve sentido incluso cuando una publicación no se
centra en Ousman Umar.

La secuencia que debería sostener la comunicación es:

\begin{quote}
Propósito → aportaciones y alianzas → recursos y formación → resultados
educativos documentados → oportunidades de desarrollo.
\end{quote}

Cada eslabón necesita evidencia. Haber abierto un aula demuestra
actividad; acreditar aprendizaje exige indicadores educativos; demostrar
un efecto sobre empleo o migración requiere una evaluación adicional. No
son resultados intercambiables.

\hypertarget{coherencia-del-relato}{%
\subsubsection{3.3. Coherencia del relato}\label{coherencia-del-relato}}

La publicación de LinkedIn del 02/09/2026 explica que la transformación
digital requiere educación básica, alfabetización digital y formación
del profesorado. Este mensaje refuerza la idea de que el valor de NASCO
no consiste únicamente en proporcionar ordenadores. Las publicaciones de
mayo de 2026 sobre aulas y entidades colaboradoras materializan ese
relato en proyectos concretos.
\href{https://www.linkedin.com/posts/nasco-feeding-minds_nascofoundation-educaci\%C3\%B3ndigital-impactosocial-activity-7500847528586846208-5q6E}{F12}
\href{https://www.linkedin.com/posts/nasco-feeding-minds_nascofoundation-feedingminds-digitaleducation-activity-7465326634426843136-PKBW}{F13}
\href{https://www.linkedin.com/posts/nasco-feeding-minds_nasco-x-deloitte-activity-7464593630444306432-8i9r}{F14}

La oportunidad aquí es equilibrar el componente emocional de la causa
con información verificable. Los mensajes de urgencia pueden
sensibilizar, pero no sustituyen a la explicación del trabajo ni a la
rendición de cuentas.

\begin{center}\rule{0.5\linewidth}{0.5pt}\end{center}

\hypertarget{propuesta-de-valor-y-producto-ampliado}{%
\subsection{4. Propuesta de valor y Producto
Ampliado}\label{propuesta-de-valor-y-producto-ampliado}}

\hypertarget{propuesta-de-valor}{%
\subsubsection{4.1. Propuesta de valor}\label{propuesta-de-valor}}

A partir de lo que he visto en la web y la memoria, la propuesta de
valor de NASCO se podría resumir así: ofrece a particulares y empresas
una forma concreta de contribuir a la educación digital en Ghana, ya sea
con dinero o con equipos, y publica información sobre lo que hace con
esas aportaciones.

Esta propuesta tiene cuatro pilares. Veamos cada uno con su evidencia y
sus límites:

\begin{longtable}[]{@{}
  >{\raggedright\arraybackslash}p{(\columnwidth - 4\tabcolsep) * \real{0.3333}}
  >{\raggedright\arraybackslash}p{(\columnwidth - 4\tabcolsep) * \real{0.3333}}
  >{\raggedright\arraybackslash}p{(\columnwidth - 4\tabcolsep) * \real{0.3333}}@{}}
\toprule\noalign{}
\begin{minipage}[b]{\linewidth}\raggedright
Pilar
\end{minipage} & \begin{minipage}[b]{\linewidth}\raggedright
Evidencia
\end{minipage} & \begin{minipage}[b]{\linewidth}\raggedright
Qué aporta y dónde falla
\end{minipage} \\
\midrule\noalign{}
\endhead
\bottomrule\noalign{}
\endlastfoot
\textbf{Sabe de educación y tecnología} & Aulas y formación descritas en
la memoria y la web.
\href{https://nascoict.org/wp-content/uploads/MEMORIA-2024.pdf}{F03}
\href{https://nascoict.org/aulas-informaticas/}{F04} & Facilita explicar
la causa / No demuestra que sea la única que hace esto. \\
\textbf{Trabaja en Ghana, en el origen} & Historia, misión y
localización de proyectos.
\href{https://nascoict.org/quienes-somos/}{F02}
\href{https://nascoict.org/wp-content/uploads/MEMORIA-2024.pdf}{F03} &
Da concreción al propósito / No prueba por sí sola un efecto sobre la
migración. \\
\textbf{Se puede colaborar de varias formas} & Donaciones, alta de
socios y contacto para entregar equipos.
\href{https://nascoict.org/aulas-informaticas/}{F04}
\href{https://nascoict.org/hazte-socio/}{F05} & Permite distintas formas
de colaboración / Falta evaluar si son fáciles de utilizar. \\
\textbf{Publica lo que hace} & Memorias y documentos de transparencia.
\href{https://nascoict.org/transparencia/}{F06} & Puede reducir
incertidumbre / Las discrepancias entre páginas debilitan su
presentación. \\
\end{longtable}

\hypertarget{oferta-principal-elegida-para-aplicar-kotler}{%
\subsubsection{4.2. Oferta principal elegida para aplicar
Kotler}\label{oferta-principal-elegida-para-aplicar-kotler}}

Para aplicar el modelo de Producto Ampliado de Kotler, me centro en
\textbf{la colaboración económica de un particular con los programas de
NASCO}. Lo que vive el donante es distinto del servicio educativo que
recibe el estudiante en Ghana, así que los separo.

\textbf{Beneficio básico (hipótesis a validar).} La idea es que el
donante contribuye a una causa que le parece significativa y siente que
su aportación sirve para algo. No lo presento como algo ya demostrado:
el estudio sectorial de la AEFr identifica el deseo de ``cambiar las
cosas desde mi posición'' como una motivación relevante, pero eso es
general y no sustituye a investigar a los colaboradores reales de NASCO.
\href{https://www.aefundraising.org/perfil-donante-2024-que-impulsa-a-los-espanoles-a-colaborar/}{F19}

\textbf{Producto real.} Una aportación puntual o periódica, con importe
y condiciones definidos. La web ofrece tres importes orientativos para
donaciones puntuales (\textbf{30€, 50€ y 95€}), una opción personalizada
y un formulario de socios con periodicidad mensual, trimestral o anual.
Las equivalencias entre dinero y actividad que aparecen en la web son
afirmaciones de NASCO; habría que contrastar con NASCO si están
actualizadas antes de usarlas en campañas.
\href{https://nascoict.org/hazte-socio/}{F05}

\textbf{Producto ampliado.} Alrededor de la aportación hay ocho
elementos que pueden aportar comodidad, información o confianza. Los he
sacado de la web, pero en la mayoría no he podido comprobar cómo
funcionan en la práctica:

\begin{longtable}[]{@{}
  >{\raggedright\arraybackslash}p{(\columnwidth - 4\tabcolsep) * \real{0.3333}}
  >{\raggedright\arraybackslash}p{(\columnwidth - 4\tabcolsep) * \real{0.3333}}
  >{\raggedright\arraybackslash}p{(\columnwidth - 4\tabcolsep) * \real{0.3333}}@{}}
\toprule\noalign{}
\begin{minipage}[b]{\linewidth}\raggedright
Elemento
\end{minipage} & \begin{minipage}[b]{\linewidth}\raggedright
Evidencia disponible
\end{minipage} & \begin{minipage}[b]{\linewidth}\raggedright
Límite de la comprobación
\end{minipage} \\
\midrule\noalign{}
\endhead
\bottomrule\noalign{}
\endlastfoot
1. Flexibilidad de colaboración & Donación puntual y aportaciones
mensuales, trimestrales o anuales.
\href{https://nascoict.org/hazte-socio/}{F05} & No he probado el cambio
de modalidad. \\
2. Facilidades de pago & Opciones de tarjeta y Bizum para donar;
formulario de socios con IBAN.
\href{https://nascoict.org/hazte-socio/}{F05} & No he completado
pagos. \\
3. Confirmación de la aportación & Se anuncia un correo de confirmación
para la donación puntual. \href{https://nascoict.org/hazte-socio/}{F05}
& No he comprobado su recepción. \\
4. Certificado de donación & El formulario permite solicitarlo y se
explica cómo pedirlo para Bizum.
\href{https://nascoict.org/hazte-socio/}{F05} & No he verificado plazos
ni emisión. \\
5. Información fiscal & Explicaciones diferenciadas para particulares y
empresas. \href{https://nascoict.org/hazte-socio/}{F05} & No sé si las
condiciones fiscales aplican igual a todos los donantes; habría que
confirmarlo con NASCO. \\
6. Rendición de cuentas & Apartados de memorias y documentación de 2017
a 2024. \href{https://nascoict.org/transparencia/}{F06} & No he leído
cada documento; no sé si hay auditoría externa. \\
7. Comunicación posterior & Opción de suscripción al boletín.
\href{https://nascoict.org/hazte-socio/}{F05} & Desconozco frecuencia,
contenidos y recepción efectiva. \\
8. Posibilidad de baja & La portada comunica que es posible darse de
baja. \href{https://nascoict.org/}{F01} & No he probado el procedimiento
ni su facilidad. \\
\end{longtable}

\hypertarget{duxf3nde-reside-la-diferenciaciuxf3n-defendible}{%
\subsubsection{4.3. Dónde reside la diferenciación
defendible}\label{duxf3nde-reside-la-diferenciaciuxf3n-defendible}}

NASCO se distingue por la combinación de propósito, especialización
digital y proyectos concretos. Pero si solo se queda ahí, cualquier ONG
similar podría copiar el discurso. Donde de verdad podría diferenciarse
de forma más sostenible es en el \textbf{producto ampliado}: que el
donante entienda a qué va su dinero, que colaborar sea fácil y que
reciba información consistente. Actualmente la web anuncia varios de
estos elementos, pero como no he colaborado económicamente con la
organización no tengo evidencia suficiente para afirmar que la
experiencia posterior ya sea buena.

La prioridad no es añadir promesas nuevas, sino comprobar y explicar
mejor las que la organización ya comunica.

\begin{center}\rule{0.5\linewidth}{0.5pt}\end{center}

\hypertarget{diagnuxf3stico-de-presencia-digital}{%
\subsection{5. Diagnóstico de presencia
digital}\label{diagnuxf3stico-de-presencia-digital}}

\hypertarget{web-y-canales-propios}{%
\subsubsection{5.1. Web y canales propios}\label{web-y-canales-propios}}

Estos son los elementos que he revisado en la web de NASCO, con lo que
vi y lo que propongo:

\begin{longtable}[]{@{}
  >{\raggedright\arraybackslash}p{(\columnwidth - 6\tabcolsep) * \real{0.2500}}
  >{\raggedright\arraybackslash}p{(\columnwidth - 6\tabcolsep) * \real{0.2500}}
  >{\raggedright\arraybackslash}p{(\columnwidth - 6\tabcolsep) * \real{0.2500}}
  >{\raggedright\arraybackslash}p{(\columnwidth - 6\tabcolsep) * \real{0.2500}}@{}}
\toprule\noalign{}
\begin{minipage}[b]{\linewidth}\raggedright
Elemento
\end{minipage} & \begin{minipage}[b]{\linewidth}\raggedright
Qué vi
\end{minipage} & \begin{minipage}[b]{\linewidth}\raggedright
Qué pienso
\end{minipage} & \begin{minipage}[b]{\linewidth}\raggedright
Qué propongo
\end{minipage} \\
\midrule\noalign{}
\endhead
\bottomrule\noalign{}
\endlastfoot
Cifras de aulas & Portada: \textbf{11}; página de aulas: \textbf{17};
``Quiénes somos'': \textbf{22}. La portada también dice 66.000 personas
y 3.766 alumnos actuales; ``Quiénes somos'' dice 60.000 estudiantes; la
página de aulas dice 6.600 y 6.000 en la misma página.
\href{https://nascoict.org/}{F01}
\href{https://nascoict.org/quienes-somos/}{F02}
\href{https://nascoict.org/aulas-informaticas/}{F04} & Es posible que 22
sea el total acumulado desde 2012 y que 11 o 17 sean cifras de un
momento dado, pero ninguna página lo explica. & Acordar una cifra de
referencia y fechar cada indicador. \\
Denominación & La portada dice ``NASCO Foundation'' y el logo y el pie
de página dicen ``NASCO Feeding Minds''. En LinkedIn también aparece
``NASCO Foundation''. \href{https://nascoict.org/}{F01}
\href{https://nascoict.org/quienes-somos/}{F02}
\href{https://www.linkedin.com/posts/nasco-feeding-minds_nascofoundation-educaci\%C3\%B3ndigital-impactosocial-activity-7500847528586846208-5q6E}{F12}
& He buscado en la web, la memoria 2018 y perfiles públicos y no he
encontrado ninguna explicación. & Explicar la relación entre nombres y
usar siempre el mismo. \\
Donaciones & \textbf{Tres importes prefijados} más opción personalizada.
\href{https://nascoict.org/hazte-socio/}{F05} & Concretan la aportación
y su destino declarado. & Revisar las equivalencias y mostrar cuándo se
actualizaron. \\
Alta de socios & \textbf{Tres periodicidades}, datos personales y
bancarios y casilla de boletín.
\href{https://nascoict.org/hazte-socio/}{F05} & Hay un recorrido de
colaboración recurrente, pero no sé dónde se abandona. & Revisar
instrucciones y medir el abandono antes de atribuirle problemas. \\
Transparencia & \textbf{Ocho ejercicios: 2017-2024}. Se observa texto de
plantilla sin sustituir: ``Your content goes here''.
\href{https://nascoict.org/transparencia/}{F06} & Hay documentación,
pero el contenido residual da la sensación de que no se ha revisado. &
Revisar presentación, enlaces y fecha de actualización. \\
Blog enlazado & La primera página muestra \textbf{cuatro entradas},
entre septiembre de 2022 y mayo de 2023; la más reciente visible es del
\textbf{31/05/2023}. El blog está disponible en castellano, inglés y
catalán. \href{https://nascoict.org/blog/}{F07} & Parece que no se
actualiza desde hace tiempo. & Retomar el blog y hacer más
publicaciones. \\
Boletín & Se ofrece suscripción, sin datos públicos de tamaño de lista o
rendimiento. \href{https://nascoict.org/hazte-socio/}{F05} & Podría ser
un punto de continuidad con el donante, pero no sé nada de él. & Pedir
información interna antes de evaluarlo. \\
\end{longtable}

Que haya documentación publicada y un formulario no demuestra que la web
convierta bien. Sin sesiones, clics, inicios y finalizaciones no puedo
atribuir a cada elemento un efecto real sobre la captación.

\hypertarget{redes-sociales-oficiales}{%
\subsubsection{5.2. Redes sociales
oficiales}\label{redes-sociales-oficiales}}

Las cuatro cuentas están enlazadas desde la web institucional. He
excluido los perfiles personales del fundador.
\href{https://nascoict.org/}{F01}

\begin{longtable}[]{@{}
  >{\raggedright\arraybackslash}p{(\columnwidth - 6\tabcolsep) * \real{0.2500}}
  >{\raggedright\arraybackslash}p{(\columnwidth - 6\tabcolsep) * \real{0.2500}}
  >{\raggedright\arraybackslash}p{(\columnwidth - 6\tabcolsep) * \real{0.2500}}
  >{\raggedright\arraybackslash}p{(\columnwidth - 6\tabcolsep) * \real{0.2500}}@{}}
\toprule\noalign{}
\begin{minipage}[b]{\linewidth}\raggedright
Canal
\end{minipage} & \begin{minipage}[b]{\linewidth}\raggedright
Seguidores y actividad
\end{minipage} & \begin{minipage}[b]{\linewidth}\raggedright
Qué dice el perfil
\end{minipage} & \begin{minipage}[b]{\linewidth}\raggedright
Cómo lo he consultado y qué no he podido ver
\end{minipage} \\
\midrule\noalign{}
\endhead
\bottomrule\noalign{}
\endlastfoot
\textbf{Instagram - @nasco\_ict} & \textbf{34K seguidores}, \textbf{283
publicaciones}, \textbf{562 cuentas seguidas}. & Biografía centrada en
migración forzada, educación digital en Ghana, fundador y CTA
``COLABORA''. & Metadatos públicos de Instagram. ``34K'' es aproximado;
no he revisado todas las publicaciones recientes.
\href{https://www.instagram.com/nasco_ict/}{F08} \\
\textbf{Facebook - NASCO Feeding Minds} & \textbf{5.790 seguidores}. &
Descripción vinculada a Ousman Umar y educación digital para las jóvenes
generaciones. & Metadatos públicos de Facebook; no he podido ver una
secuencia reciente de publicaciones.
\href{https://www.facebook.com/NascoICT/}{F09} \\
\textbf{LinkedIn - Nasco Feeding Minds} & \textbf{2.415 seguidores}. &
Educación, proyectos y colaboraciones institucionales en la muestra que
he revisado. & Datos estructurados y publicaciones oficiales. He
verificado tres publicaciones; no tengo acceso al historial completo.
\href{https://www.linkedin.com/company/nasco-feeding-minds/}{F10}
\href{https://www.linkedin.com/posts/nasco-feeding-minds_nascofoundation-educaci\%C3\%B3ndigital-impactosocial-activity-7500847528586846208-5q6E}{F12}
\href{https://www.linkedin.com/posts/nasco-feeding-minds_nascofoundation-feedingminds-digitaleducation-activity-7465326634426843136-PKBW}{F13}
\href{https://www.linkedin.com/posts/nasco-feeding-minds_nasco-x-deloitte-activity-7464593630444306432-8i9r}{F14} \\
\textbf{X - @nasco\_ict} & \textbf{2.400 seguidores}, \textbf{1.753
publicaciones}, \textbf{790 cuentas seguidas}. & Biografía bilingüe
sobre educación digital en Ghana; la muestra visible combina noticias y
sensibilización migratoria. & Lectura pública del perfil. Las últimas
publicaciones subidas datan del 2024, por lo que la cuenta está,
aparentemente inactiva desde entonces.
\href{https://twitter.com/nasco_ict/}{F11} \\
\end{longtable}

Instagram es la cuenta con más seguidores de las que he visto. Eso no
significa que sea el canal que capta más donantes. Tampoco se pueden
sumar las cuatro cifras como si fueran una audiencia única, porque puede
haber gente que siga a NASCO en varias redes a la vez.

\hypertarget{muestra-de-contenido-en-linkedin}{%
\subsubsection{5.3. Muestra de contenido en
LinkedIn}\label{muestra-de-contenido-en-linkedin}}

He elegido tres publicaciones de LinkedIn porque me sirven para analizar
el propósito, los proyectos y las alianzas de NASCO:

\begin{longtable}[]{@{}
  >{\raggedright\arraybackslash}p{(\columnwidth - 6\tabcolsep) * \real{0.2308}}
  >{\raggedright\arraybackslash}p{(\columnwidth - 6\tabcolsep) * \real{0.2308}}
  >{\raggedleft\arraybackslash}p{(\columnwidth - 6\tabcolsep) * \real{0.3077}}
  >{\raggedright\arraybackslash}p{(\columnwidth - 6\tabcolsep) * \real{0.2308}}@{}}
\toprule\noalign{}
\begin{minipage}[b]{\linewidth}\raggedright
Fecha
\end{minipage} & \begin{minipage}[b]{\linewidth}\raggedright
Qué publican
\end{minipage} & \begin{minipage}[b]{\linewidth}\raggedleft
Reacciones
\end{minipage} & \begin{minipage}[b]{\linewidth}\raggedright
Qué me dice
\end{minipage} \\
\midrule\noalign{}
\endhead
\bottomrule\noalign{}
\endlastfoot
\textbf{02/09/2026} & Educación básica, alfabetización digital y
formación del profesorado. Texto en español e imágenes.
\href{https://www.linkedin.com/posts/nasco-feeding-minds_nascofoundation-educaci\%C3\%B3ndigital-impactosocial-activity-7500847528586846208-5q6E}{F12}
& \textbf{4} & Explica por qué entregar tecnología no es suficiente y
desarrolla el CÓMO de la marca. \\
\textbf{27/05/2026} & Aula de educación digital en Tuna con Barça
Foundation. Texto en catalán e inglés; vídeo de \textbf{49 segundos}.
\href{https://www.linkedin.com/posts/nasco-feeding-minds_nascofoundation-feedingminds-digitaleducation-activity-7465326634426843136-PKBW}{F13}
& \textbf{29} & Hace tangible un proyecto y reconoce a la entidad
colaboradora. \\
\textbf{25/05/2026} & Centro de formación digital con Deloitte en Wa
Islamic Girls Senior High School. Texto en español e inglés.
\href{https://www.linkedin.com/posts/nasco-feeding-minds_nasco-x-deloitte-activity-7464593630444306432-8i9r}{F14}
& \textbf{26} & Relaciona colaboración empresarial, formación y
oportunidades para jóvenes estudiantes. \\
\end{longtable}

\hypertarget{tono-consistencia-y-diferenciaciuxf3n-digital}{%
\subsubsection{5.4. Tono, consistencia y diferenciación
digital}\label{tono-consistencia-y-diferenciaciuxf3n-digital}}

\textbf{Fortaleza: se nota de qué va.} Las descripciones y piezas que he
revisado mantienen educación digital, Ghana y oportunidades como ejes.
La muestra de LinkedIn permite ver actividad más allá de la historia del
fundador.

\textbf{Tono: emocional e institucional.} Los mensajes de esperanza
conviven con explicaciones y colaboraciones. Esta combinación puede
acompañar a personas con distintos niveles de conocimiento de la causa,
siempre que las afirmaciones de impacto estén respaldadas.

\textbf{Comunicación multilingüe.} La presencia de español, catalán e
inglés responde a un contexto internacional. Para captar en España
convendría fijar un criterio que facilite la comprensión, sin asumir que
todos los usuarios entienden las tres lenguas.

\textbf{Debilidad: los datos no cuadran.} Los nombres y cifras no
siempre se presentan de forma uniforme. En una organización que solicita
confianza, esta inconsistencia afecta a la claridad del relato, aunque
no haya medido su efecto sobre las donaciones.

\textbf{Límite del diagnóstico.} Analizo lo que NASCO proyecta
públicamente, no la percepción real de todos sus públicos. Para conocer
reputación y confianza sería necesario incorporar entrevistas,
comentarios representativos o investigación adicional.

\begin{center}\rule{0.5\linewidth}{0.5pt}\end{center}

\hypertarget{anuxe1lisis-del-entorno-pestel}{%
\subsection{6. Análisis del entorno:
PESTEL}\label{anuxe1lisis-del-entorno-pestel}}

Aquí separo lo que es externo de lo que es problema interno de NASCO.
Por ejemplo, la desconfianza hacia las ONG es un factor social; que las
cifras de aulas no cuadren en la web es un problema de la propia marca.
Lo que propongo en cada factor son ideas, no cosas que NASCO haya hecho.

\hypertarget{poluxedtico}{%
\subsubsection{6.1. Político}\label{poluxedtico}}

\begin{longtable}[]{@{}
  >{\raggedright\arraybackslash}p{(\columnwidth - 8\tabcolsep) * \real{0.2000}}
  >{\raggedright\arraybackslash}p{(\columnwidth - 8\tabcolsep) * \real{0.2000}}
  >{\raggedright\arraybackslash}p{(\columnwidth - 8\tabcolsep) * \real{0.2000}}
  >{\raggedright\arraybackslash}p{(\columnwidth - 8\tabcolsep) * \real{0.2000}}
  >{\raggedright\arraybackslash}p{(\columnwidth - 8\tabcolsep) * \real{0.2000}}@{}}
\toprule\noalign{}
\begin{minipage}[b]{\linewidth}\raggedright
Código y factor
\end{minipage} & \begin{minipage}[b]{\linewidth}\raggedright
Evidencia y periodo
\end{minipage} & \begin{minipage}[b]{\linewidth}\raggedright
Clasificación
\end{minipage} & \begin{minipage}[b]{\linewidth}\raggedright
Implicación para NASCO
\end{minipage} & \begin{minipage}[b]{\linewidth}\raggedright
Acción concreta propuesta
\end{minipage} \\
\midrule\noalign{}
\endhead
\bottomrule\noalign{}
\endlastfoot
\textbf{P1. Contracción de la ayuda oficial al desarrollo} & Los datos
preliminares de la OCDE para \textbf{2025}, recogidos por la Comisión
Europea en \textbf{2026}, indican una caída real del \textbf{23,1 \%}
respecto a 2024.
\href{https://knowledge4policy.ec.europa.eu/publication/oecd-preliminary-official-development-assistance-levels-2025_en}{F15}
& Amenaza sectorial potencial. & Puede presionar las fuentes de
financiación del sector, aunque no significa que los ingresos de NASCO
vayan a bajar lo mismo. & Preparar \textbf{una pieza} que explique por
qué seguir colaborando sostiene los programas, sin atribuir a NASCO
recortes que no están confirmados. \\
\textbf{P2. Marco español favorable a la cooperación} & La \textbf{Ley
1/2023}, publicada el \textbf{21/02/2023}, fija el compromiso de
destinar al menos el \textbf{0,7 \% de la RNB a AOD en 2030}.
\href{https://www.boe.es/buscar/act.php?id=BOE-A-2023-4512}{F16} &
Oportunidad de alineamiento institucional. & Hay un contexto favorable a
la cooperación, pero no garantiza subvenciones ni que Ghana sea
prioridad en una convocatoria concreta. & Hacer \textbf{una ficha} que
conecte la actividad educativa con desarrollo y colaboración, usando
solo resultados acreditados. \\
\end{longtable}

\hypertarget{econuxf3mico}{%
\subsubsection{6.2. Económico}\label{econuxf3mico}}

\begin{longtable}[]{@{}
  >{\raggedright\arraybackslash}p{(\columnwidth - 8\tabcolsep) * \real{0.2000}}
  >{\raggedright\arraybackslash}p{(\columnwidth - 8\tabcolsep) * \real{0.2000}}
  >{\raggedright\arraybackslash}p{(\columnwidth - 8\tabcolsep) * \real{0.2000}}
  >{\raggedright\arraybackslash}p{(\columnwidth - 8\tabcolsep) * \real{0.2000}}
  >{\raggedright\arraybackslash}p{(\columnwidth - 8\tabcolsep) * \real{0.2000}}@{}}
\toprule\noalign{}
\begin{minipage}[b]{\linewidth}\raggedright
Código y factor
\end{minipage} & \begin{minipage}[b]{\linewidth}\raggedright
Evidencia y periodo
\end{minipage} & \begin{minipage}[b]{\linewidth}\raggedright
Clasificación
\end{minipage} & \begin{minipage}[b]{\linewidth}\raggedright
Implicación para NASCO
\end{minipage} & \begin{minipage}[b]{\linewidth}\raggedright
Acción concreta propuesta
\end{minipage} \\
\midrule\noalign{}
\endhead
\bottomrule\noalign{}
\endlastfoot
\textbf{E1. Restricciones económicas de los hogares españoles} & La ECV
\textbf{2025}, publicada por el INE el \textbf{05/02/2026}, registra un
\textbf{36,4 \%} de población sin capacidad para afrontar gastos
imprevistos. \href{https://www.ine.es/dyngs/Prensa/es/ECV2025.htm}{F17}
& Amenaza para el compromiso económico. & Que a alguien le interese una
causa no significa que pueda mantener una cuota. & Mostrar
\textbf{juntas} la donación puntual, la colaboración periódica y otras
vías, con condiciones claras y sin presión. \\
\textbf{E2. Costes logísticos y riesgos económicos en Ghana} & El Banco
Mundial, el \textbf{26/08/2026}, señala recuperación económica, pero
también riesgos energéticos y limitaciones del transporte. Más del
\textbf{95 \%} del tráfico de pasajeros y mercancías se realiza por
carretera.
\href{https://www.worldbank.org/en/news/press-release/2026/08/26/ghana-must-turn-strong-recovery-into-quality-jobs-and-lasting-growth}{F18}
& Amenaza operativa potencial. & Llevar y mantener equipos depende de
costes e infraestructuras que una campaña no controla. & Revisar
\textbf{cada equivalencia monetaria} que se publique y hacer \textbf{una
explicación} del mantenimiento y la logística, validada con la
organización. \\
\end{longtable}

\hypertarget{social}{%
\subsubsection{6.3. Social}\label{social}}

\begin{longtable}[]{@{}
  >{\raggedright\arraybackslash}p{(\columnwidth - 8\tabcolsep) * \real{0.2000}}
  >{\raggedright\arraybackslash}p{(\columnwidth - 8\tabcolsep) * \real{0.2000}}
  >{\raggedright\arraybackslash}p{(\columnwidth - 8\tabcolsep) * \real{0.2000}}
  >{\raggedright\arraybackslash}p{(\columnwidth - 8\tabcolsep) * \real{0.2000}}
  >{\raggedright\arraybackslash}p{(\columnwidth - 8\tabcolsep) * \real{0.2000}}@{}}
\toprule\noalign{}
\begin{minipage}[b]{\linewidth}\raggedright
Código y factor
\end{minipage} & \begin{minipage}[b]{\linewidth}\raggedright
Evidencia y periodo
\end{minipage} & \begin{minipage}[b]{\linewidth}\raggedright
Clasificación
\end{minipage} & \begin{minipage}[b]{\linewidth}\raggedright
Implicación para NASCO
\end{minipage} & \begin{minipage}[b]{\linewidth}\raggedright
Acción concreta propuesta
\end{minipage} \\
\midrule\noalign{}
\endhead
\bottomrule\noalign{}
\endlastfoot
\textbf{S1. Desconfianza en la gestión de las ONG} & La AEFr, en su
publicación del \textbf{04/12/2024}, recoge la desconfianza en gestión y
transparencia como barrera para el \textbf{20 \% de los no donantes}.
\href{https://www.aefundraising.org/perfil-donante-2024-que-impulsa-a-los-espanoles-a-colaborar/}{F19}
& Amenaza; oportunidad de responder con evidencia. & Si las cifras, la
actividad y la rendición de cuentas no son coherentes, la desconfianza
aumenta. & Crear \textbf{una página de resultados} con periodo,
definición y fuente para cada indicador. \\
\textbf{S2. Motivación para contribuir a un cambio concreto} & La misma
publicación identifica ``cambiar las cosas desde mi posición'' como
principal motivación y recoge que el \textbf{38,2 \%} de la población de
referencia colaboró durante el último año. Estudio \textbf{2024}.
\href{https://www.aefundraising.org/perfil-donante-2024-que-impulsa-a-los-espanoles-a-colaborar/}{F19}
& Oportunidad. & Hay gente dispuesta a colaborar, pero eso no significa
que NASCO vaya a captarlos ni permite estimar conversiones. &
Desarrollar \textbf{un caso educativo documentado} que conecte
aportación, actividad y resultado, con autorización para usar
testimonios e imágenes. \\
\end{longtable}

\hypertarget{tecnoluxf3gico}{%
\subsubsection{6.4. Tecnológico}\label{tecnoluxf3gico}}

\begin{longtable}[]{@{}
  >{\raggedright\arraybackslash}p{(\columnwidth - 8\tabcolsep) * \real{0.2000}}
  >{\raggedright\arraybackslash}p{(\columnwidth - 8\tabcolsep) * \real{0.2000}}
  >{\raggedright\arraybackslash}p{(\columnwidth - 8\tabcolsep) * \real{0.2000}}
  >{\raggedright\arraybackslash}p{(\columnwidth - 8\tabcolsep) * \real{0.2000}}
  >{\raggedright\arraybackslash}p{(\columnwidth - 8\tabcolsep) * \real{0.2000}}@{}}
\toprule\noalign{}
\begin{minipage}[b]{\linewidth}\raggedright
Código y factor
\end{minipage} & \begin{minipage}[b]{\linewidth}\raggedright
Evidencia y periodo
\end{minipage} & \begin{minipage}[b]{\linewidth}\raggedright
Clasificación
\end{minipage} & \begin{minipage}[b]{\linewidth}\raggedright
Implicación para NASCO
\end{minipage} & \begin{minipage}[b]{\linewidth}\raggedright
Acción concreta propuesta
\end{minipage} \\
\midrule\noalign{}
\endhead
\bottomrule\noalign{}
\endlastfoot
\textbf{T1. Uso extendido de Internet en España} & El INE señala que en
\textbf{2025} el \textbf{96,3 \%} de las personas de 16-74 años utilizó
Internet en los tres meses anteriores. Publicación: \textbf{20/11/2025}.
\href{https://www.ine.es/dyngs/Prensa/es/TICH2025.htm}{F20} &
Oportunidad de acceso digital. & Justifica trabajar la web y los
contenidos, pero no demuestra que la gente quiera donar ni prefiera una
red concreta. & Hacer una prueba del recorrido de colaboración con
\textbf{cinco personas}, sin pagos reales, para ver dónde dudan o se
atascan. \\
\textbf{T2. Expansión del uso de IA generativa} & En la misma encuesta,
el \textbf{37,9 \%} de las personas de 16-74 años utilizó IA generativa
en los tres meses anteriores.
\href{https://www.ine.es/dyngs/Prensa/es/TICH2025.htm}{F20} &
Oportunidad y amenaza. & Puede ayudar a producir contenidos, pero
también hace más importante que lo que se publique sea exacto y esté
supervisado. & Crear \textbf{una ficha de revisión} para cada contenido
hecho con IA: fuente, fecha, afirmaciones, revisión y aprobación. \\
\end{longtable}

\hypertarget{ecoluxf3gico}{%
\subsubsection{6.5. Ecológico}\label{ecoluxf3gico}}

\begin{longtable}[]{@{}
  >{\raggedright\arraybackslash}p{(\columnwidth - 8\tabcolsep) * \real{0.2000}}
  >{\raggedright\arraybackslash}p{(\columnwidth - 8\tabcolsep) * \real{0.2000}}
  >{\raggedright\arraybackslash}p{(\columnwidth - 8\tabcolsep) * \real{0.2000}}
  >{\raggedright\arraybackslash}p{(\columnwidth - 8\tabcolsep) * \real{0.2000}}
  >{\raggedright\arraybackslash}p{(\columnwidth - 8\tabcolsep) * \real{0.2000}}@{}}
\toprule\noalign{}
\begin{minipage}[b]{\linewidth}\raggedright
Código y factor
\end{minipage} & \begin{minipage}[b]{\linewidth}\raggedright
Evidencia y periodo
\end{minipage} & \begin{minipage}[b]{\linewidth}\raggedright
Clasificación
\end{minipage} & \begin{minipage}[b]{\linewidth}\raggedright
Implicación para NASCO
\end{minipage} & \begin{minipage}[b]{\linewidth}\raggedright
Acción concreta propuesta
\end{minipage} \\
\midrule\noalign{}
\endhead
\bottomrule\noalign{}
\endlastfoot
\textbf{EC1. Residuos electrónicos y necesidad de reutilización
responsable} & El Global E-waste Monitor, publicado el
\textbf{20/03/2024}, registra \textbf{62 millones de toneladas} de
residuos electrónicos en \textbf{2022}; solo el \textbf{22,3 \%} se
documentó como recogido y reciclado formalmente.
\href{https://www.itu.int/en/ITU-D/Environment/Pages/Publications/The-Global-E-waste-Monitor-2024.aspx}{F21}
& Oportunidad y amenaza. & Reutilizar equipos conecta colaboración
empresarial y educación, pero si se envía material inservible se
traslada el problema. & Preparar \textbf{una guía de equipos aceptados},
estado requerido y qué se hace con ellos al final; validar el
procedimiento antes de prometerlo. \\
\textbf{EC2. Vulnerabilidad climática de Ghana} & El Banco Mundial, el
\textbf{01/11/2022}, describe daños potenciales en infraestructuras por
lluvias erráticas e inundaciones y estima unas \textbf{45.000 personas
afectadas anualmente}, según aquel diagnóstico.
\href{https://www.worldbank.org/en/news/press-release/2022/11/01/ghana-can-turn-climate-challenges-into-opportunities-for-resilient-and-sustainable-growth-says-new-world-bank-group-report}{F22}
& Amenaza operativa. & Mantener la educación funcionando puede requerir
reparaciones y adaptaciones. No afirmo que una instalación de NASCO haya
sufrido daños ni que la cifra sea de 2026. & Añadir \textbf{una
explicación de mantenimiento y continuidad} a la información de
proyectos. \\
\end{longtable}

\hypertarget{legal}{%
\subsubsection{6.6. Legal}\label{legal}}

\begin{longtable}[]{@{}
  >{\raggedright\arraybackslash}p{(\columnwidth - 8\tabcolsep) * \real{0.2000}}
  >{\raggedright\arraybackslash}p{(\columnwidth - 8\tabcolsep) * \real{0.2000}}
  >{\raggedright\arraybackslash}p{(\columnwidth - 8\tabcolsep) * \real{0.2000}}
  >{\raggedright\arraybackslash}p{(\columnwidth - 8\tabcolsep) * \real{0.2000}}
  >{\raggedright\arraybackslash}p{(\columnwidth - 8\tabcolsep) * \real{0.2000}}@{}}
\toprule\noalign{}
\begin{minipage}[b]{\linewidth}\raggedright
Código y factor
\end{minipage} & \begin{minipage}[b]{\linewidth}\raggedright
Evidencia y periodo
\end{minipage} & \begin{minipage}[b]{\linewidth}\raggedright
Clasificación
\end{minipage} & \begin{minipage}[b]{\linewidth}\raggedright
Implicación para NASCO
\end{minipage} & \begin{minipage}[b]{\linewidth}\raggedright
Acción concreta propuesta
\end{minipage} \\
\midrule\noalign{}
\endhead
\bottomrule\noalign{}
\endlastfoot
\textbf{L1. Condiciones para el consentimiento de cookies} & La AEPD, el
\textbf{11/07/2023}, indicó que aceptar y rechazar deben situarse al
mismo nivel, sin dificultar el rechazo. Adaptación exigida como máximo
el \textbf{11/01/2024}.
\href{https://www.aepd.es/prensa-y-comunicacion/notas-de-prensa/aepd-actualiza-guia-cookies-para-adaptarla-a-nuevas-directrices-cepd}{F23}
& Riesgo de incumplimiento; oportunidad de confianza. & Si se quiere
medir, hay que respetar el consentimiento. La revisión que he hecho no
es una auditoría legal de NASCO. & Revisar \textbf{el banner, la
configuración y la activación de analítica} antes de lanzar campañas, y
anotar las incidencias que se detecten. \\
\textbf{L2. Incentivos fiscales sujetos a requisitos} & El Manual de
Renta \textbf{2024} de la AEAT, actualizado el \textbf{10/04/2025},
recoge el \textbf{80 \% sobre los primeros 250 €}, el \textbf{40 \%
sobre el resto} y un tipo incrementado del \textbf{45 \%} sujeto a
reiteración y condiciones.
\href{https://sede.agenciatributaria.gob.es/Sede/ayuda/manuales-videos-folletos/manuales-practicos/irpf-2024/c16-deducciones-generales-cuota/deducciones-donativos-otras-aportaciones/donativos-realizados-entidades-incluidas/general-donativos-donaciones-entidades-beneficiarias-mecenazgo/porcentajes-deduccion.html}{F24}
& Oportunidad informativa; riesgo de promesas incorrectas. & Puede
ayudar a entender la colaboración, pero no se puede anunciar como un
ahorro para todo el mundo. Hay que revisar ejercicio, requisitos,
límites y territorio. & Hacer \textbf{una FAQ fiscal} con fuente y
fecha; confirmar que aplica a NASCO y al ejercicio correspondiente antes
de publicarla. \\
\end{longtable}

\hypertarget{cinco-factores-externos-prioritarios}{%
\subsubsection{6.7. Cinco factores externos
prioritarios}\label{cinco-factores-externos-prioritarios}}

He ordenado estos cinco factores según mi criterio. He tenido en cuenta
qué tan cerca están de la decisión de colaborar, qué graves pueden ser y
qué puede hacer NASCO desde contenidos o experiencia. Sin datos internos
no se puede hacer un cálculo riguroso.

\begin{longtable}[]{@{}
  >{\raggedleft\arraybackslash}p{(\columnwidth - 4\tabcolsep) * \real{0.4000}}
  >{\raggedright\arraybackslash}p{(\columnwidth - 4\tabcolsep) * \real{0.3000}}
  >{\raggedright\arraybackslash}p{(\columnwidth - 4\tabcolsep) * \real{0.3000}}@{}}
\toprule\noalign{}
\begin{minipage}[b]{\linewidth}\raggedleft
Prioridad
\end{minipage} & \begin{minipage}[b]{\linewidth}\raggedright
Factor
\end{minipage} & \begin{minipage}[b]{\linewidth}\raggedright
Justificación
\end{minipage} \\
\midrule\noalign{}
\endhead
\bottomrule\noalign{}
\endlastfoot
1 & \textbf{S1. Confianza} & Se relaciona directamente con la
información que la marca debe ofrecer antes de pedir una aportación. \\
2 & \textbf{E1. Capacidad económica} & Condiciona el tipo de compromiso
que puede asumir un particular. \\
3 & \textbf{T1. Acceso digital} & Sustenta la posibilidad de atraer y
acompañar al público a través de canales propios. \\
4 & \textbf{L1. Consentimiento} & Condiciona cómo puede implantarse la
medición digital propuesta. \\
5 & \textbf{EC1. Reutilización responsable} & Es especialmente relevante
para la colaboración B2B mediante equipos. \\
\end{longtable}

\begin{center}\rule{0.5\linewidth}{0.5pt}\end{center}

\hypertarget{tendencias-relevantes-consumer-trends-2026}{%
\subsection{7. Tendencias relevantes: Consumer Trends
2026}\label{tendencias-relevantes-consumer-trends-2026}}

He mirado la presentación pública de Consumer Trends 2026 de Zorraquino
y las páginas de tres tendencias. No me he descargado el informe
completo, el que se pide por formulario. Las tendencias me sirven como
ideas generales, no como datos de los donantes de NASCO.
\href{https://trends.zorraquino.com/informes/consumer-trends/2026/}{F25}

He elegido estas tres porque tienen que ver con educación y captación
digital, encajan con la marca y se pueden convertir en contenidos. No
cojo tendencias de otros sectores sin justificar por qué tienen que ver
con la causa y el público de NASCO.

\begin{longtable}[]{@{}
  >{\raggedright\arraybackslash}p{(\columnwidth - 6\tabcolsep) * \real{0.2500}}
  >{\raggedright\arraybackslash}p{(\columnwidth - 6\tabcolsep) * \real{0.2500}}
  >{\raggedright\arraybackslash}p{(\columnwidth - 6\tabcolsep) * \real{0.2500}}
  >{\raggedright\arraybackslash}p{(\columnwidth - 6\tabcolsep) * \real{0.2500}}@{}}
\toprule\noalign{}
\begin{minipage}[b]{\linewidth}\raggedright
Tendencia y fuente
\end{minipage} & \begin{minipage}[b]{\linewidth}\raggedright
Lectura relevante
\end{minipage} & \begin{minipage}[b]{\linewidth}\raggedright
Oportunidad o riesgo para NASCO
\end{minipage} & \begin{minipage}[b]{\linewidth}\raggedright
Aplicación propuesta y comprobación
\end{minipage} \\
\midrule\noalign{}
\endhead
\bottomrule\noalign{}
\endlastfoot
\textbf{Brain rot - tendencia 6.}
\href{https://trends.zorraquino.com/informes/consumer-trends/2026/brain-rot.html}{F26}
& Zorraquino analiza la saturación y superficialidad del consumo
digital. No adopto su etiqueta como diagnóstico clínico. & Oportunidad
de distinguir educación digital de exposición pasiva a pantallas; riesgo
de añadir contenido repetitivo sin utilidad. & Producir \textbf{tres
piezas educativas} sobre qué se aprende en un aula, cómo trabaja el
docente y cómo se comprueba el aprendizaje. Valorar comprensión con
\textbf{cinco lectores}, sin fijar de antemano resultados. \\
\textbf{Cognitive co-agents - tendencia 9.}
\href{https://trends.zorraquino.com/informes/consumer-trends/2026/cognitiveco-agents.html}{F27}
& La IA participa cada vez más en procesos educativos y culturales, con
riesgos de homogeneización y desigualdad de acceso. & Oportunidad para
explicar el valor de las competencias digitales; riesgo de prometer
formación en IA que NASCO no haya acreditado ofrecer. & Elaborar
\textbf{una pieza} sobre competencias, criterio y acompañamiento humano.
Someter todas las afirmaciones sobre programas de NASCO a validación
previa. \\
\textbf{Everything campaignable - tendencia 4.}
\href{https://trends.zorraquino.com/informes/consumer-trends/2026/everythingcampaignable.html}{F28}
& La presión por convertir acontecimientos en campañas puede aumentar la
visibilidad, pero diluir la identidad o parecer oportunista. &
Oportunidad de aprovechar hitos educativos y colaboraciones; riesgo de
utilizar la migración o las historias personales solo para generar
atención. & Crear \textbf{un criterio editorial de cinco preguntas}:
relación con la misión, evidencia, consentimiento, dignidad de las
personas y siguiente paso útil para el lector. \\
\end{longtable}

Estas aplicaciones no exigen abrir nuevas redes ni aumentar la
frecuencia sin criterio. Favorecen contenidos coherentes y verificables,
que es una necesidad que ya detecté en el análisis de marca.

\begin{center}\rule{0.5\linewidth}{0.5pt}\end{center}

\hypertarget{suxedntesis-estratuxe9gica-del-diagnuxf3stico-a-las-prioridades}{%
\subsection{8. Síntesis estratégica: del diagnóstico a las
prioridades}\label{suxedntesis-estratuxe9gica-del-diagnuxf3stico-a-las-prioridades}}

\hypertarget{tres-oportunidades-externas}{%
\subsubsection{8.1. Tres oportunidades
externas}\label{tres-oportunidades-externas}}

\begin{enumerate}
\def\labelenumi{\arabic{enumi}.}
\tightlist
\item
  \textbf{Público digital accesible:} el uso extendido de Internet
  permite desarrollar recorridos de información y colaboración, sin
  depender únicamente de acciones presenciales.
  \href{https://www.ine.es/dyngs/Prensa/es/TICH2025.htm}{F20}
\item
  \textbf{Motivación para contribuir:} la AEFr identifica el deseo de
  generar cambios; NASCO puede concretarlo mediante proyectos educativos
  y evidencia, no solo con mensajes emocionales.
  \href{https://www.aefundraising.org/perfil-donante-2024-que-impulsa-a-los-espanoles-a-colaborar/}{F19}
\item
  \textbf{Reutilización tecnológica responsable:} el problema de los
  residuos electrónicos permite explicar una colaboración empresarial
  que conecte recursos y educación, siempre que se documenten
  condiciones y destino.
  \href{https://www.itu.int/en/ITU-D/Environment/Pages/Publications/The-Global-E-waste-Monitor-2024.aspx}{F21}
\end{enumerate}

\hypertarget{tres-amenazas-externas}{%
\subsubsection{8.2. Tres amenazas
externas}\label{tres-amenazas-externas}}

\begin{enumerate}
\def\labelenumi{\arabic{enumi}.}
\tightlist
\item
  \textbf{Desconfianza en la gestión:} obliga a sostener el relato con
  información clara y consistente.
  \href{https://www.aefundraising.org/perfil-donante-2024-que-impulsa-a-los-espanoles-a-colaborar/}{F19}
\item
  \textbf{Limitaciones económicas y presión sobre la cooperación:}
  pueden dificultar los compromisos recurrentes; no permiten anticipar
  por sí solas una caída de ingresos en NASCO.
  \href{https://knowledge4policy.ec.europa.eu/publication/oecd-preliminary-official-development-assistance-levels-2025_en}{F15}
  \href{https://www.ine.es/dyngs/Prensa/es/ECV2025.htm}{F17}
\item
  \textbf{Saturación informativa y riesgo de mensajes oportunistas:} la
  búsqueda de visibilidad puede perjudicar la identidad y la
  credibilidad si sustituye a la utilidad.
  \href{https://trends.zorraquino.com/informes/consumer-trends/2026/brain-rot.html}{F26}
  \href{https://trends.zorraquino.com/informes/consumer-trends/2026/everythingcampaignable.html}{F28}
\end{enumerate}

\hypertarget{cinco-prioridades-propuestas}{%
\subsubsection{8.3. Cinco prioridades
propuestas}\label{cinco-prioridades-propuestas}}

\begin{longtable}[]{@{}
  >{\raggedleft\arraybackslash}p{(\columnwidth - 4\tabcolsep) * \real{0.4000}}
  >{\raggedright\arraybackslash}p{(\columnwidth - 4\tabcolsep) * \real{0.3000}}
  >{\raggedright\arraybackslash}p{(\columnwidth - 4\tabcolsep) * \real{0.3000}}@{}}
\toprule\noalign{}
\begin{minipage}[b]{\linewidth}\raggedleft
Orden
\end{minipage} & \begin{minipage}[b]{\linewidth}\raggedright
Prioridad y justificación
\end{minipage} & \begin{minipage}[b]{\linewidth}\raggedright
Resultado que permitiría comprobar el avance
\end{minipage} \\
\midrule\noalign{}
\endhead
\bottomrule\noalign{}
\endlastfoot
1 & \textbf{Llegar a más personas en España mostrando primero resultados
educativos documentados.} Es la prioridad que he elegido para este
trabajo: tanto al particular como a la empresa les mostraría un caso
concreto antes de pedirles colaboración. Los seguidores no indican por
sí solos a cuántos colaboradores nuevos se llega.
\href{https://nascoict.org/wp-content/uploads/MEMORIA-2024.pdf}{F03}
\href{https://www.instagram.com/nasco_ict/}{F08}
\href{https://www.facebook.com/NascoICT/}{F09} & Publicar un caso con
periodo y fuente, y separar la procedencia de visitas y solicitudes de
particulares y empresas; comparar periodos sin equiparar seguidores con
donantes. \\
2 & \textbf{Explicar los resultados a quien descubre la causa.} Hay
memorias públicas y la transparencia importa para los potenciales
donantes.
\href{https://nascoict.org/wp-content/uploads/MEMORIA-2024.pdf}{F03}
\href{https://nascoict.org/transparencia/}{F06}
\href{https://www.aefundraising.org/perfil-donante-2024-que-impulsa-a-los-espanoles-a-colaborar/}{F19}
& Un caso educativo fechado que separe actividad, aprendizaje e impacto,
con fuente y autorización para imágenes o testimonios. \\
3 & \textbf{Distinguir a particulares y empresas sin complicar la
donación.} El recorrido general de colaboración me resulta claro, pero
una empresa que aporta equipos puede necesitar información diferente.
\href{https://nascoict.org/aulas-informaticas/}{F04}
\href{https://nascoict.org/hazte-socio/}{F05} & Contenidos de entrada
diferenciados; conservar las vías de colaboración existentes y
contrastar qué dudas aparecen realmente. \\
4 & \textbf{Aclarar los datos contradictorios.} La portada, la página de
aulas y ``Quiénes somos'' no muestran la misma cifra.
\href{https://nascoict.org/}{F01}
\href{https://nascoict.org/quienes-somos/}{F02}
\href{https://nascoict.org/aulas-informaticas/}{F04} & Fechar y definir
cada indicador; comprobar si las diferencias obedecen a periodos o tipos
de aula distintos. \\
5 & \textbf{Medir y revisar antes de escalar.} No hay acceso a
analítica, CRM ni presupuesto de marketing. & Establecer una línea de
base y responsables; revisar el consentimiento antes de atribuir altas a
las campañas.
\href{https://www.aepd.es/prensa-y-comunicacion/notas-de-prensa/aepd-actualiza-guia-cookies-para-adaptarla-a-nuevas-directrices-cepd}{F23} \\
\end{longtable}

El primer puesto recoge una decisión mía. Los demás son propuestas
argumentadas, no compromisos operativos de NASCO; el orden puede cambiar
cuando se disponga de datos internos.

\hypertarget{conclusiuxf3n-del-diagnuxf3stico}{%
\subsubsection{8.4. Conclusión del
diagnóstico}\label{conclusiuxf3n-del-diagnuxf3stico}}

En resumen, NASCO tiene proyectos documentados y presencia en redes,
pero no sé si está llegando a gente nueva en España. Por eso he puesto
como primera prioridad que más personas descubran la organización, tanto
particulares como empresas. No creo que el problema esté en la página de
colaboración: la he revisado y me parece clara. El problema es que la
gente no llega hasta ella. Y cuando llega, las cifras que se encuentra
no cuadran, lo que no ayuda a generar confianza. Más adelante propongo
objetivos y una forma de medirlos, dejando claro qué datos me faltan.

\begin{center}\rule{0.5\linewidth}{0.5pt}\end{center}

\hypertarget{recursos-humanos-y-capacidad-de-ejecuciuxf3n}{%
\subsection{9. Recursos humanos y capacidad de
ejecución}\label{recursos-humanos-y-capacidad-de-ejecuciuxf3n}}

La memoria de 2024 nombra a Ousman Umar y Banasco Seidu como fundadores
y dice que han hecho cambios para estabilizar la financiación. La página
de aulas menciona profesorado en Ghana. Pero con lo que hay público no
puedo saber cuánta gente trabaja en marketing ni cuánto tiempo le
dedican.
\href{https://nascoict.org/wp-content/uploads/MEMORIA-2024.pdf}{F03}
\href{https://nascoict.org/aulas-informaticas/}{F04}

Para la estrategia que propongo haría falta cubrir estas funciones:

\begin{longtable}[]{@{}
  >{\raggedright\arraybackslash}p{(\columnwidth - 4\tabcolsep) * \real{0.3333}}
  >{\raggedright\arraybackslash}p{(\columnwidth - 4\tabcolsep) * \real{0.3333}}
  >{\raggedright\arraybackslash}p{(\columnwidth - 4\tabcolsep) * \real{0.3333}}@{}}
\toprule\noalign{}
\begin{minipage}[b]{\linewidth}\raggedright
Función
\end{minipage} & \begin{minipage}[b]{\linewidth}\raggedright
Qué hay hoy
\end{minipage} & \begin{minipage}[b]{\linewidth}\raggedright
Qué habría que preguntar a NASCO
\end{minipage} \\
\midrule\noalign{}
\endhead
\bottomrule\noalign{}
\endlastfoot
Coordinar mensajes y revisar datos de impacto & Hay web, memorias y
publicaciones. \href{https://nascoict.org/}{F01}
\href{https://nascoict.org/wp-content/uploads/MEMORIA-2024.pdf}{F03} &
Quién valida las cifras educativas y autoriza su uso. \\
Publicar y atender a posibles colaboradores & Hay redes oficiales y
canales de contacto. \href{https://nascoict.org/}{F01}
\href{https://nascoict.org/hazte-socio/}{F05} & Quién lo hace, cuánto
tiempo tiene y tarda en responder. \\
Mantener web, formularios y medición & Hay páginas de colaboración, pero
no tengo acceso a su configuración.
\href{https://nascoict.org/hazte-socio/}{F05} & Quién gestiona el sitio,
si hay analítica y cómo se recoge el consentimiento. \\
Preparar contenidos sobre los programas de Ghana & Las memorias y
publicaciones muestran proyectos.
\href{https://nascoict.org/wp-content/uploads/MEMORIA-2024.pdf}{F03}
\href{https://www.linkedin.com/posts/nasco-feeding-minds_nascofoundation-educaci\%C3\%B3ndigital-impactosocial-activity-7500847528586846208-5q6E}{F12}
\href{https://www.linkedin.com/posts/nasco-feeding-minds_nascofoundation-feedingminds-digitaleducation-activity-7465326634426843136-PKBW}{F13}
& Quién consigue materiales del terreno y permisos para testimonios e
imágenes. \\
\end{longtable}

No voy a inventar un organigrama ni a asignar tareas a personas que no
conozco. Antes de recomendar si esto se hace por dentro o se
externaliza, hace falta saber si NASCO tiene gente para ello.

\hypertarget{recursos-econuxf3micos}{%
\subsection{10. Recursos económicos}\label{recursos-econuxf3micos}}

No tengo el presupuesto de marketing de NASCO. La web menciona
\textbf{15.000 € anuales para el mantenimiento de cada aula}, pero
\textbf{eso no es presupuesto publicitario}. Tampoco puedo sacar de la
memoria cuánto se gasta en comunicación.
\href{https://nascoict.org/}{F01}
\href{https://nascoict.org/wp-content/uploads/MEMORIA-2024.pdf}{F03}

Para ver si una estrategia modesta sería viable, propongo este escenario
de seis meses con \textbf{1.200 €}. Es un supuesto para el trabajo, no
dinero que NASCO vaya a gastarse:

\begin{longtable}[]{@{}
  >{\raggedright\arraybackslash}p{(\columnwidth - 4\tabcolsep) * \real{0.3000}}
  >{\raggedright\arraybackslash}p{(\columnwidth - 4\tabcolsep) * \real{0.3000}}
  >{\raggedleft\arraybackslash}p{(\columnwidth - 4\tabcolsep) * \real{0.4000}}@{}}
\toprule\noalign{}
\begin{minipage}[b]{\linewidth}\raggedright
Partida hipotética
\end{minipage} & \begin{minipage}[b]{\linewidth}\raggedright
Supuesto de cálculo
\end{minipage} & \begin{minipage}[b]{\linewidth}\raggedleft
Coste de referencia
\end{minipage} \\
\midrule\noalign{}
\endhead
\bottomrule\noalign{}
\endlastfoot
Distribución pagada de contenido de captación & 100 €/mes durante seis
meses, \textbf{solo si se autoriza}. & 600 € \\
Adaptación puntual de piezas y página de colaboración & Reserva de
trabajo externo, si no existe capacidad interna. & 300 € \\
Contingencia y pruebas & Reserva para correcciones o pequeñas pruebas
autorizadas. & 300 € \\
\textbf{Total supuesto} & \textbf{600 + 300 + 300} & \textbf{1.200 €} \\
\end{longtable}

Esto no incluye horas del equipo, equipos informáticos ni lo que cuesta
operar en Ghana; esos costes no los conozco. Si NASCO no tiene esta
cantidad, la alternativa es publicar de forma orgánica con las
herramientas que ya tiene, siempre que haya alguien para hacerlo.
\textbf{Antes de usar las metas de la sección 14 hay que confirmar
presupuesto, gente disponible y datos de partida}.

\hypertarget{competencia-y-alternativas-para-la-colaboraciuxf3n}{%
\subsection{11. Competencia y alternativas para la
colaboración}\label{competencia-y-alternativas-para-la-colaboraciuxf3n}}

Cuando aquí digo ``competencia'' me refiero a \textbf{otras opciones por
las que un particular o una empresa podría apoyar educación o aportar
tecnología}, no a organizaciones enemigas de NASCO. He elegido tres:
TeSo, Entreculturas y Worldreader. La comparación es cualitativa,
mirando sus páginas web. No he comparado resultados, costes de captación
ni cuota de mercado.

\begin{longtable}[]{@{}
  >{\raggedright\arraybackslash}p{(\columnwidth - 6\tabcolsep) * \real{0.2500}}
  >{\raggedright\arraybackslash}p{(\columnwidth - 6\tabcolsep) * \real{0.2500}}
  >{\raggedright\arraybackslash}p{(\columnwidth - 6\tabcolsep) * \real{0.2500}}
  >{\raggedright\arraybackslash}p{(\columnwidth - 6\tabcolsep) * \real{0.2500}}@{}}
\toprule\noalign{}
\begin{minipage}[b]{\linewidth}\raggedright
Entidad
\end{minipage} & \begin{minipage}[b]{\linewidth}\raggedright
Qué hace y qué dice
\end{minipage} & \begin{minipage}[b]{\linewidth}\raggedright
Qué comunica bien
\end{minipage} & \begin{minipage}[b]{\linewidth}\raggedright
Dónde se diferencia de NASCO y qué se puede aprender
\end{minipage} \\
\midrule\noalign{}
\endhead
\bottomrule\noalign{}
\endlastfoot
\textbf{TeSo}, asociación española de reutilización tecnológica. &
Explica qué equipos acepta y cómo preparar una donación; documenta la
entrega de \textbf{12 equipos} para un aula escolar en Ghana.
\href{https://teso.org.es/informacion/donar-material/}{F31}
\href{https://teso.org.es/entrega-del-proyecto-oyera-d-a-school-de-ghana/}{F32}
& Condiciones técnicas y destino del material concretos. & Se centra en
reutilizar y ceder equipos a entidades, no en un modelo de socios para
sostener aulas. NASCO podría explicar mejor los requisitos para donar
equipos. \\
\textbf{Entreculturas}, alternativa para donantes de educación en
España. & Presenta cuotas para particulares, un recorrido específico
para empresas y preguntas frecuentes sobre la colaboración.
\href{https://www.entreculturas.org/hazte-socio/}{F33}
\href{https://www.entreculturas.org/hazte-socio-empresas/}{F34} &
Distingue ambos públicos y responde dudas después del alta. & Trabaja en
muchos países y causas; NASCO se especializa en educación digital en
Ghana, que es otra propuesta, no una ventaja de resultados
demostrada. \\
\textbf{Worldreader}, referente sectorial internacional. & Comunica
lectura digital y colaboración con familias; documentó un programa en
Ghana con Compassion International y ofrece vías de donación.
\href{https://www.worldreader.org/ways-to-donate/}{F35}
\href{https://www.worldreader.org/now/worldreader-and-compassion-international-join-forces/}{F36}
& Relaciona tecnología y hábito educativo, con contenidos para familias.
& Es una organización internacional con vías fiscales y de donación
diferentes; no la trato como competidor directo en España. Sirve de
referencia para explicar aprendizaje, no solo acceso al dispositivo. \\
\end{longtable}

No incluyo a ProFuturo como competidor: su web documenta una alianza
educativa con NASCO en Ghana, así que es un aliado.
\href{https://profuturo.education/noticias/profuturo-nasco-feeding-minds-educacion-digital-ghana/}{F37}
Lo que digo de las otras entidades es solo en comparación con el
objetivo de este trabajo, no una evaluación de lo bien o mal que lo
hacen por dentro. Estas tres alternativas representan lo que un donante
o empresa podría considerar antes de elegir NASCO, y se conectan con el
diagnóstico de marca (sección 5) y con los buyer personas (sección 12).

\hypertarget{buyer-persona-hipuxf3tesis-de-trabajo}{%
\subsection{12. Buyer persona: hipótesis de
trabajo}\label{buyer-persona-hipuxf3tesis-de-trabajo}}

No he entrevistado a donantes de NASCO ni he visto su CRM, así que estos
dos perfiles son hipótesis. Los he construido a partir del análisis de
marca y del estudio sectorial de la AEFr, que me da una idea general de
qué motiva a la gente a colaborar.
\href{https://www.aefundraising.org/perfil-donante-2024-que-impulsa-a-los-espanoles-a-colaborar/}{F19}

\begin{longtable}[]{@{}
  >{\raggedright\arraybackslash}p{(\columnwidth - 4\tabcolsep) * \real{0.3333}}
  >{\raggedright\arraybackslash}p{(\columnwidth - 4\tabcolsep) * \real{0.3333}}
  >{\raggedright\arraybackslash}p{(\columnwidth - 4\tabcolsep) * \real{0.3333}}@{}}
\toprule\noalign{}
\begin{minipage}[b]{\linewidth}\raggedright
Aspecto
\end{minipage} & \begin{minipage}[b]{\linewidth}\raggedright
B2C: posible socio particular en España
\end{minipage} & \begin{minipage}[b]{\linewidth}\raggedright
B2B: responsable de colaboración en una empresa española
\end{minipage} \\
\midrule\noalign{}
\endhead
\bottomrule\noalign{}
\endlastfoot
Necesidad & Entender qué hace NASCO antes de comprometerse todos los
meses. & Encontrar un proyecto que encaje con los valores y recursos de
su empresa. \\
Información que mostraría primero & Un caso de resultados educativos con
periodo y fuente, antes de explicar la cuota. & Resultados documentados
que pueda enseñar a otras personas de la empresa, antes de hablar de
dinero. \\
Motivación \textbf{hipotética} & Contribuir a oportunidades educativas
concretas y recibir noticias sobre lo conseguido. & Dar utilidad social
a una aportación económica o a equipos aptos, con resultados que pueda
comunicar. \\
Dolor u objeción & Duda sobre a dónde va el dinero o dificultad para
elegir una causa entre muchas. & Incertidumbre sobre qué material sirve,
cómo se lleva, qué hay que aprobar internamente y cómo se hace
seguimiento. \\
Criterio de decisión & Proyecto que entienda, cifras con fecha, opciones
de cuota y confianza en la entidad. & Condiciones de participación,
encaje social, documentación y persona de contacto. \\
Fuentes que \textbf{podría mirar} & Web de NASCO, buscadores,
publicaciones y memorias; habría que confirmarlo con personas reales. &
LinkedIn, web, memoria y vías de contacto; habría que confirmarlo con
empresas reales. \\
Conversión propuesta & Alta de socio con primera aportación confirmada.
& Solicitud de colaboración identificada y cualificada; un contacto
todavía no es un convenio. \\
\end{longtable}

Si pudiera hablar con colaboradores reales, podría confirmar qué
objeciones de la tabla importan de verdad, pero los perfiles se
sostienen a partir del diagnóstico que he hecho. El recorrido completo
de cada perfil está en el customer journey (sección 13).

\hypertarget{customer-journey}{%
\subsection{13. Customer journey}\label{customer-journey}}

El recorrido no describe un comportamiento ya observado: muestra
\textbf{qué habría que acompañar y medir}. Las etapas de particulares y
empresas comparten lógica, pero no la misma conversión. Cada etapa se
vincula al buyer persona correspondiente (sección 12) y los KPIs remiten
a los objetivos SMART (sección 14).

\begin{longtable}[]{@{}
  >{\raggedright\arraybackslash}p{(\columnwidth - 6\tabcolsep) * \real{0.2500}}
  >{\raggedright\arraybackslash}p{(\columnwidth - 6\tabcolsep) * \real{0.2500}}
  >{\raggedright\arraybackslash}p{(\columnwidth - 6\tabcolsep) * \real{0.2500}}
  >{\raggedright\arraybackslash}p{(\columnwidth - 6\tabcolsep) * \real{0.2500}}@{}}
\toprule\noalign{}
\begin{minipage}[b]{\linewidth}\raggedright
Etapa y éxito esperado
\end{minipage} & \begin{minipage}[b]{\linewidth}\raggedright
Particular: contacto, necesidad y freno
\end{minipage} & \begin{minipage}[b]{\linewidth}\raggedright
Empresa: contacto, necesidad y freno
\end{minipage} & \begin{minipage}[b]{\linewidth}\raggedright
Contenido y medida propuesta
\end{minipage} \\
\midrule\noalign{}
\endhead
\bottomrule\noalign{}
\endlastfoot
\textbf{Descubrimiento:} conoce NASCO. & Llega por búsqueda o red; entre
muchas causas, necesita un ejemplo que le muestre qué se ha hecho. &
Encuentra un caso o publicación; necesita resultados que pueda explicar
dentro de su empresa. & Presentar primero un caso educativo fechado y
con fuente. Medir sesiones de España con origen identificado y clics
hacia esa página. \\
\textbf{Consideración:} entiende la actividad. & Contrasta el caso con
la memoria; necesita distinguir actividad realizada de aprendizaje
demostrado. & Revisa documentación del proyecto antes de valorar
aportaciones y condiciones. & Resumen de resultados y límites, seguido
de preguntas frecuentes y vías de colaboración; medir visitas a
proyectos y a la página de colaboración, sin equipararlas a intención de
donar. \\
\textbf{Decisión:} solicita o aporta. & Compara donación puntual con
alta periódica; el freno sería una duda real, \textbf{no una dificultad
de uso que aquí no he comprobado}. & Contacta con NASCO; puede necesitar
concretar destino, requisitos o logística. & Vías de colaboración
existentes y contacto; medir primera aportación confirmada o solicitud
empresarial cualificada. \\
\textbf{Vinculación:} mantiene la relación. & Quiere saber cómo continúa
el proyecto; posible freno: poca información posterior, pendiente de
probar. & Necesita seguimiento para justificar la relación ante su
equipo; pendiente de confirmar. & Actualización de actividad y respuesta
a dudas; medir cuotas efectivamente renovadas y colaboraciones
activas. \\
\textbf{Recomendación:} comparte con otros. & Recomienda la causa si
confía en lo vivido. & Puede difundir el proyecto con autorización y
datos contrastados. & Historia verificable y posibilidad de recomendar;
medir referencias identificadas sin suponer que una mención pública
genera donantes. \\
\end{longtable}

\hypertarget{bruxfajula-estratuxe9gica-y-objetivos-smart}{%
\subsection{14. Brújula estratégica y objetivos
SMART}\label{bruxfajula-estratuxe9gica-y-objetivos-smart}}

\textbf{Síntesis de diagnóstico:} existen cuatro comunidades sociales
oficiales y vías de colaboración visibles, pero no se conoce su
contribución real a nuevas relaciones. Considero clara la página de
colaboración y priorizo \textbf{llegar a más personas}. Son riesgos la
saturación de mensajes, las dudas habituales sobre transparencia y las
cifras institucionales sin periodo homogéneo. Las \textbf{cinco
prioridades máximas} se ordenan en la sección 8.3.
\href{https://nascoict.org/}{F01}
\href{https://nascoict.org/transparencia/}{F06}
\href{https://www.instagram.com/nasco_ict/}{F08}
\href{https://www.facebook.com/NascoICT/}{F09}
\href{https://www.aefundraising.org/perfil-donante-2024-que-impulsa-a-los-espanoles-a-colaborar/}{F19}

\textbf{Objetivo de negocio propuesto:} ampliar en España una base
sostenible de personas socias y empresas colaboradoras para apoyar los
programas educativos. La aspiración de 2.000 socios mencionada en la
memoria de 2024 \textbf{no es la base actual ni la traslado directamente
como objetivo de esta campaña}.
\href{https://nascoict.org/wp-content/uploads/MEMORIA-2024.pdf}{F03}

\begin{longtable}[]{@{}
  >{\raggedright\arraybackslash}p{(\columnwidth - 6\tabcolsep) * \real{0.2500}}
  >{\raggedright\arraybackslash}p{(\columnwidth - 6\tabcolsep) * \real{0.2500}}
  >{\raggedright\arraybackslash}p{(\columnwidth - 6\tabcolsep) * \real{0.2500}}
  >{\raggedright\arraybackslash}p{(\columnwidth - 6\tabcolsep) * \real{0.2500}}@{}}
\toprule\noalign{}
\begin{minipage}[b]{\linewidth}\raggedright
Objetivo propuesto
\end{minipage} & \begin{minipage}[b]{\linewidth}\raggedright
Criterio de consecución y plazo
\end{minipage} & \begin{minipage}[b]{\linewidth}\raggedright
Relación con el diagnóstico
\end{minipage} & \begin{minipage}[b]{\linewidth}\raggedright
Dato que falta antes de aprobarlo
\end{minipage} \\
\midrule\noalign{}
\endhead
\bottomrule\noalign{}
\endlastfoot
\textbf{O1. Atracción:} aumentar un \textbf{25 \%} las sesiones
identificadas desde España que llegan a contenidos educativos de NASCO.
& Comparar los últimos 30 días del \textbf{sexto mes} de campaña con los
30 días de línea de base anteriores al comienzo; excluir tráfico interno
y automatizado si la herramienta lo permite. & Comprobar si se llega a
más personas; los seguidores actuales no miden ese avance. & Valor de
partida, definición de sesión y disponibilidad de analítica con
consentimiento. El 25 \% es una \textbf{meta de prueba}, no validada con
NASCO. \\
\textbf{O2. B2C:} conseguir \textbf{20 nuevas altas de socios
particulares} procedentes de canales digitales en \textbf{seis meses},
con primera aportación registrada. & Contar altas únicas de residentes
en España confirmadas en el periodo; no contabilizar intentos o visitas.
& Traducir parte del nuevo conocimiento de la marca en apoyo recurrente.
& Base actual, capacidad de seguimiento y validación del umbral de
20. \\
\textbf{O3. B2B:} registrar \textbf{seis solicitudes empresariales
cualificadas} de colaboración en \textbf{seis meses}. & Empresa
identificada, contacto, tipo de aportación propuesto y procedencia
digital; no equivale a seis acuerdos. & Trabajar el segundo público sin
confundir interés con ingresos confirmados. & Histórico de solicitudes,
capacidad de respuesta y validación del umbral de seis. \\
\end{longtable}

He decidido mantener las tres metas: +25 \% de sesiones, 20 nuevas altas
y seis solicitudes empresariales en seis meses. Tienen formulación
específica, criterio de medida y plazo. Sin datos internos de NASCO no
puedo saber si son alcanzables, pero me parecen razonables y se
ajustarían con la organización antes de lanzar una campaña. El horizonte
empieza cuando haya una fecha real de puesta en marcha.

Cada objetivo se mide mediante los KPIs que desarrollo en la sección 15:
\textbf{O1} se mide con sesiones desde España y clics orgánicos
(atracción); \textbf{O2} se mide con altas confirmadas y retención de
cohorte (conversión y fidelización); \textbf{O3} se mide con solicitudes
cualificadas y colaboraciones continuadas (conversión y fidelización).

\hypertarget{kpis-y-plan-de-mediciuxf3n}{%
\subsection{15. KPIs y plan de
medición}\label{kpis-y-plan-de-mediciuxf3n}}

Una métrica pública de seguidores sirve para describir el punto de
partida, pero no sustituye a medir visitas, colaboraciones y
permanencia. Las herramientas de la tabla son propuestas; no sé si NASCO
tiene GA4, Search Console, UTMs o un CRM. Si se usan, hay que respetar
los requisitos de consentimiento.
\href{https://www.aepd.es/prensa-y-comunicacion/notas-de-prensa/aepd-actualiza-guia-cookies-para-adaptarla-a-nuevas-directrices-cepd}{F23}
Los baselines son estimaciones, no datos reales de NASCO; los
responsables son roles que se confirmarían con la organización.

\begin{longtable}[]{@{}
  >{\raggedright\arraybackslash}p{(\columnwidth - 12\tabcolsep) * \real{0.1429}}
  >{\raggedright\arraybackslash}p{(\columnwidth - 12\tabcolsep) * \real{0.1429}}
  >{\raggedright\arraybackslash}p{(\columnwidth - 12\tabcolsep) * \real{0.1429}}
  >{\raggedright\arraybackslash}p{(\columnwidth - 12\tabcolsep) * \real{0.1429}}
  >{\raggedright\arraybackslash}p{(\columnwidth - 12\tabcolsep) * \real{0.1429}}
  >{\raggedright\arraybackslash}p{(\columnwidth - 12\tabcolsep) * \real{0.1429}}
  >{\raggedright\arraybackslash}p{(\columnwidth - 12\tabcolsep) * \real{0.1429}}@{}}
\toprule\noalign{}
\begin{minipage}[b]{\linewidth}\raggedright
KPI vinculado a
\end{minipage} & \begin{minipage}[b]{\linewidth}\raggedright
Etapa del embudo
\end{minipage} & \begin{minipage}[b]{\linewidth}\raggedright
KPI y definición operativa
\end{minipage} & \begin{minipage}[b]{\linewidth}\raggedright
Baseline estimado
\end{minipage} & \begin{minipage}[b]{\linewidth}\raggedright
Target
\end{minipage} & \begin{minipage}[b]{\linewidth}\raggedright
Fuente y periodicidad
\end{minipage} & \begin{minipage}[b]{\linewidth}\raggedright
Responsable propuesto
\end{minipage} \\
\midrule\noalign{}
\endhead
\bottomrule\noalign{}
\endlastfoot
\textbf{O1} & Atracción & Sesiones desde España a páginas educativas con
origen atribuible. Crecimiento =
\texttt{(sesiones\ fin\ -\ sesiones\ base)\ /\ sesiones\ base\ ×\ 100}.
& \textbf{300 sesiones/mes} (supuesto de partida para una web de una ONG
pequeña con \textasciitilde5.800 seguidores totales en Facebook). &
\textbf{+25 \% → 375 sesiones/mes} al sexto mes. & GA4 previa revisión
de consentimiento y UTMs; \textbf{mensual}. & Responsable digital (rol
por confirmar con NASCO). \\
\textbf{O1} & Atracción & Clics orgánicos hacia contenidos relevantes y
búsquedas que los generan; indicador de descubrimiento, no donaciones. &
\textbf{80 clics/mes} (supuesto de partida). & Crecimiento proporcional
al objetivo de sesiones. & Search Console \textbf{si existe acceso};
\textbf{mensual}. & Responsable digital (rol por confirmar). \\
\textbf{O2} & Conversión & Nuevas altas únicas de socios de España con
\textbf{primera aportación registrada} y canal identificable; contar
personas, no formularios iniciados. & \textbf{3 altas/mes} (supuesto de
partida sin datos internos). & \textbf{20 altas en seis meses}
(≈3,3/mes, alineado con el baseline). & Registro de socios y origen de
campaña, \textbf{mensual}. & Responsable de socios (rol por
confirmar). \\
\textbf{O3} & Conversión & Solicitudes B2B cualificadas: empresa
identificada + contacto + interés concreto + origen, deduplicadas. &
\textbf{1 solicitud/mes} (supuesto de partida). & \textbf{6 solicitudes
en seis meses} (≈1/mes, alineado con el baseline). & Registro de
contactos y correo/formulario, \textbf{mensual}. & Responsable de
alianzas (rol por confirmar). \\
\textbf{O2} & Fidelización & Retención de la cohorte de nuevos socios al
tercer mes =
\texttt{socios\ activos\ al\ mes\ 3\ /\ socios\ con\ primera\ cuota\ confirmada\ ×\ 100}.
& \textbf{70 \%} (estimación basada en la media sectorial de AEFr para
pequeñas ONG; dato orientativo). & \textbf{≥75 \%} al tercer mes de cada
cohorte. & Registro de cuotas, \textbf{trimestral}. & Responsable de
socios (rol por confirmar). \\
\textbf{O3} & Fidelización & Empresas con colaboración continuada o
acuerdo activo al cierre del periodo / empresas con primera colaboración
confirmada × 100. & \textbf{Sin baseline} (no hay acuerdos previos
conocidos). & Informar del número, no de un porcentaje, hasta que
existan acuerdos. & Registro de alianzas, \textbf{trimestral}. &
Responsable de alianzas (rol por confirmar). \\
\end{longtable}

\textbf{Secuencia para hacer operativo el plan:} 1) confirmar
responsables y permiso de uso de datos con NASCO; 2) comprobar
consentimiento y registrar al menos 30 días de partida real; 3) acordar
definición de cada evento con la entidad; 4) validar presupuesto y
metas; 5) revisar mensualmente resultados sin atribuir causalidad a un
canal por una simple visita.

\hypertarget{luxedmites-y-validaciones-pendientes}{%
\subsection{16. Límites y validaciones
pendientes}\label{luxedmites-y-validaciones-pendientes}}

\begin{longtable}[]{@{}
  >{\raggedright\arraybackslash}p{(\columnwidth - 4\tabcolsep) * \real{0.3333}}
  >{\raggedright\arraybackslash}p{(\columnwidth - 4\tabcolsep) * \real{0.3333}}
  >{\raggedright\arraybackslash}p{(\columnwidth - 4\tabcolsep) * \real{0.3333}}@{}}
\toprule\noalign{}
\begin{minipage}[b]{\linewidth}\raggedright
Cuestión
\end{minipage} & \begin{minipage}[b]{\linewidth}\raggedright
Qué se conoce
\end{minipage} & \begin{minipage}[b]{\linewidth}\raggedright
Qué no se debe afirmar todavía
\end{minipage} \\
\midrule\noalign{}
\endhead
\bottomrule\noalign{}
\endlastfoot
Número actual de aulas y alumnos & La memoria de 2024 declara 17 aulas y
6.000 alumnos; la portada dice 11 aulas, 66.000 personas y 3.766 alumnos
actuales; ``Quiénes somos'' dice 22 aulas y 60.000 estudiantes; la
página de aulas dice 6.600 y 6.000 en la misma página. & Que cualquiera
de ellas sea la cifra vigente en septiembre de 2026 sin confirmación. Es
posible que algunas sean acumuladas y otras de un momento dado, pero no
se explica. \\
Número de socios & La memoria plantea llegar a 2.000. & Que la
organización tenga actualmente 2.000 socios. \\
Financiación pública & La versión inglesa afirma que no recibe
subvenciones gubernamentales; la memoria de 2024 describe participación
en REFFECT AFRICA, financiado por la Comisión Europea.
\href{https://nascoict.org/wp-content/uploads/MEMORIA-2024.pdf}{F03}
(pp.~9-10) \href{https://nascoict.org/en/}{F30} & Que exista
contradicción contable demostrada o, en sentido contrario, que no haya
ninguna relación con financiación pública. Hay que aclarar financiación
directa, participación y periodo. \\
Rendimiento de redes & Contadores y tres publicaciones de LinkedIn
verificadas. & Frecuencia mensual, crecimiento, alcance, engagement
representativo o donaciones atribuibles. \\
Servicio al colaborador & Condiciones y elementos anunciados en la web.
& Que confirmaciones, certificados, atención y bajas funcionen con una
calidad o plazo concretos. \\
Impacto & Actividad educativa declarada y proyectos comunicados. & Que
abrir aulas demuestre por sí solo creación de empleo o reducción de
migración. \\
Marco legal y fiscal & Fuentes institucionales y mensajes públicos. &
Cumplimiento integral o deducciones garantizadas para todos los
donantes. \\
\end{longtable}

Las metas y el presupuesto son supuestos que se ajustarían con NASCO
antes de una ejecución real.

\begin{center}\rule{0.5\linewidth}{0.5pt}\end{center}

\hypertarget{registro-de-decisiones-e-iteraciones}{%
\subsection{Registro de decisiones e
iteraciones}\label{registro-de-decisiones-e-iteraciones}}

\hypertarget{feedback-recibido-y-cuxf3mo-lo-he-integrado}{%
\subsubsection{Feedback recibido y cómo lo he
integrado}\label{feedback-recibido-y-cuxf3mo-lo-he-integrado}}

La docente me dijo que el documento estaba bien estructurado y tenía
fuentes citadas, pero que no era mío: se notaba que no había tocado una
coma, y el propio anexo de IA avisaba de que la revisión humana no se
había hecho. Me pidió que lo leyera, comprobara las fuentes,
reescribiera con mi criterio, cambiara lo que no diría y decidiera yo
las prioridades. Esto es lo que he hecho:

\begin{longtable}[]{@{}
  >{\raggedright\arraybackslash}p{(\columnwidth - 4\tabcolsep) * \real{0.3333}}
  >{\raggedright\arraybackslash}p{(\columnwidth - 4\tabcolsep) * \real{0.3333}}
  >{\raggedright\arraybackslash}p{(\columnwidth - 4\tabcolsep) * \real{0.3333}}@{}}
\toprule\noalign{}
\begin{minipage}[b]{\linewidth}\raggedright
Qué me dijeron
\end{minipage} & \begin{minipage}[b]{\linewidth}\raggedright
Qué he hecho
\end{minipage} & \begin{minipage}[b]{\linewidth}\raggedright
Impacto esperado
\end{minipage} \\
\midrule\noalign{}
\endhead
\bottomrule\noalign{}
\endlastfoot
``El documento no es tuyo, y se nota que no has tocado una coma.'' & He
leído cada apartado, comprobado las fuentes y reescrito el dossier con
mis palabras y mis decisiones. & El dossier refleja mi criterio, no el
de la IA, y puedo defenderlo. \\
``Estás aquí para aprender a pensar una estrategia, no para delegar ese
aprendizaje en una herramienta.'' & He sido yo quien ha decidido las
prioridades, las metas, los buyer persona y qué mostrar primero a cada
público. & He aprendido a razonar la estrategia, no solo a
describirla. \\
``La IA es un acelerador de tu criterio, no un sustituto.'' & He usado
la IA para buscar fuentes y redactar borradores, pero las decisiones y
la revisión son mías. & El anexo de trazabilidad deja claro qué hizo la
IA y qué hice yo. \\
``Cambia lo que no dirías con tus palabras.'' & He reescrito el lenguaje
para que suene a mí, no a un informe de consultoría. & El dossier se lee
como un trabajo propio, no generado. \\
\end{longtable}

\hypertarget{decisiones-propias}{%
\subsubsection{Decisiones propias}\label{decisiones-propias}}

\begin{longtable}[]{@{}
  >{\raggedright\arraybackslash}p{(\columnwidth - 4\tabcolsep) * \real{0.3333}}
  >{\raggedright\arraybackslash}p{(\columnwidth - 4\tabcolsep) * \real{0.3333}}
  >{\raggedright\arraybackslash}p{(\columnwidth - 4\tabcolsep) * \real{0.3333}}@{}}
\toprule\noalign{}
\begin{minipage}[b]{\linewidth}\raggedright
Decisión o cambio
\end{minipage} & \begin{minipage}[b]{\linewidth}\raggedright
Motivo
\end{minipage} & \begin{minipage}[b]{\linewidth}\raggedright
Consecuencia en esta versión
\end{minipage} \\
\midrule\noalign{}
\endhead
\bottomrule\noalign{}
\endlastfoot
Revisar las fuentes del borrador & He comprobado personalmente las
cifras de aulas, los seguidores en redes, la página de colaboración y la
memoria 2024. & Se conservan las fuentes previas y las nuevas F31-F37,
que también he revisado. \\
No priorizar una reforma del recorrido de donación & Tras mirarlo de
nuevo, me parece claro cómo colaborar. & Se deja de afirmar que la
donación es el problema principal; las eventuales fricciones quedan como
hipótesis, no como hecho. \\
Priorizar llegar a más personas en España & Es la estrategia que elijo
para este trabajo. & Cambian resumen, prioridades y brújula; se separan
particulares y empresas y se propone medir procedencia de visitas y
contactos. \\
Ampliar al EA1 final & La entrega final incluye cliente, recursos,
competencia, objetivos y medición. & Se añaden las secciones 9-15 con mi
revisión y mis decisiones. \\
Mantener metas, escenario económico y referentes & He decidido conservar
+25 \% de sesiones, 20 altas, seis solicitudes B2B, 1.200 € en seis
meses y la comparación con TeSo, Entreculturas y Worldreader. & Se
explicita que son supuestos académicos, no resultados observados ni
compromisos de NASCO. \\
Empezar por resultados educativos para ambos públicos & Tanto a
particulares como a empresas les mostraría primero resultados
documentados antes de pedirles colaboración. & Se ajustan las
prioridades, los dos buyer persona y el comienzo del customer
journey. \\
\end{longtable}

\begin{center}\rule{0.5\linewidth}{0.5pt}\end{center}

\hypertarget{anexo.-trazabilidad-del-uso-de-ia}{%
\subsection{Anexo. Trazabilidad del uso de
IA}\label{anexo.-trazabilidad-del-uso-de-ia}}

\hypertarget{herramientas-finalidad-y-alcance}{%
\subsubsection{1. Herramientas, finalidad y
alcance}\label{herramientas-finalidad-y-alcance}}

\textbf{Herramienta de IA:} Devin, utilizado como corrector de
estructura del documento y para la preparación de borradores.\\
\textbf{Funciones utilizadas:} búsqueda web, lectura de páginas
públicas, consulta de metadatos y datos estructurados, propuesta de
análisis y redacción de borradores.\\
\textbf{Herramienta auxiliar:} Python con \texttt{pypdf} para extraer
texto de los PDF; no es una herramienta generativa.

Devin me ayudó a estructurar el documento y a redactar borradores. Yo he
revisado todas las secciones, comprobado las fuentes, corregido una
prioridad estratégica y tomado las decisiones que constan en el registro
de iteraciones.

\hypertarget{quuxe9-verifiquuxe9-yo-personalmente}{%
\subsubsection{2. Qué verifiqué yo
personalmente}\label{quuxe9-verifiquuxe9-yo-personalmente}}

\begin{longtable}[]{@{}
  >{\raggedright\arraybackslash}p{(\columnwidth - 4\tabcolsep) * \real{0.3333}}
  >{\raggedright\arraybackslash}p{(\columnwidth - 4\tabcolsep) * \real{0.3333}}
  >{\raggedright\arraybackslash}p{(\columnwidth - 4\tabcolsep) * \real{0.3333}}@{}}
\toprule\noalign{}
\begin{minipage}[b]{\linewidth}\raggedright
Fuente o dato
\end{minipage} & \begin{minipage}[b]{\linewidth}\raggedright
Qué comprobé
\end{minipage} & \begin{minipage}[b]{\linewidth}\raggedright
Conclusión que saqué
\end{minipage} \\
\midrule\noalign{}
\endhead
\bottomrule\noalign{}
\endlastfoot
\textbf{Cifras de aulas} \href{https://nascoict.org/}{F01}
\href{https://nascoict.org/quienes-somos/}{F02}
\href{https://nascoict.org/aulas-informaticas/}{F04} & Entré en la
portada, la página de aulas y ``Quiénes somos'' y vi que muestran 11, 17
y 22 sin fechar de forma homogénea. & Es un problema de consistencia que
debe aclararse antes de reutilizar esas cifras en nuevos contenidos. \\
\textbf{Seguidores en redes}
\href{https://www.instagram.com/nasco_ict/}{F08}
\href{https://www.facebook.com/NascoICT/}{F09}
\href{https://twitter.com/nasco_ict/}{F11}
\href{https://www.linkedin.com/posts/nasco-feeding-minds_nascofoundation-educaci\%C3\%B3ndigital-impactosocial-activity-7500847528586846208-5q6E}{F12}
& Entré en Instagram, Facebook, LinkedIn y X y contrasté las cifras de
seguidores citadas. & NASCO tiene presencia en las cuatro redes, pero
los seguidores no miden alcance ni conversión. \\
\textbf{Página de colaboración}
\href{https://nascoict.org/hazte-socio/}{F05} & Revisé el formulario de
socio, los importes orientativos y las opciones de donación. & Me
pareció claro cómo colaborar; por eso no lo planteo como el problema
principal. \\
\textbf{Memoria 2024}
\href{https://nascoict.org/wp-content/uploads/MEMORIA-2024.pdf}{F03} &
Leí las páginas clave de la memoria anual 2024. & Confirma 17 aulas y
6.000 alumnos en 2024, plantea actualizar la web y aspira a 2.000
socios; no es la situación actual. \\
\end{longtable}

\hypertarget{quuxe9-produjo-la-ia-y-quuxe9-descartuxe9-o-correguxed}{%
\subsubsection{3. Qué produjo la IA y qué descarté o
corregí}\label{quuxe9-produjo-la-ia-y-quuxe9-descartuxe9-o-correguxed}}

\begin{longtable}[]{@{}
  >{\raggedright\arraybackslash}p{(\columnwidth - 2\tabcolsep) * \real{0.5000}}
  >{\raggedright\arraybackslash}p{(\columnwidth - 2\tabcolsep) * \real{0.5000}}@{}}
\toprule\noalign{}
\begin{minipage}[b]{\linewidth}\raggedright
Contenido generado por la IA
\end{minipage} & \begin{minipage}[b]{\linewidth}\raggedright
Qué hice yo
\end{minipage} \\
\midrule\noalign{}
\endhead
\bottomrule\noalign{}
\endlastfoot
Prioridad inicial: ``la página de donación es confusa''. & La descarté
tras revisarla: me parece clara. Cambié la prioridad a llegar a más
personas. \\
Advertencias repetidas sobre límites en cada sección. & Las consolidé en
un solo apartado de límites para mejorar la lectura. \\
Propuesta de metas: +25 \% sesiones, 20 altas, 6 contactos B2B, 1.200 €.
& Las revisé y decidí mantenerlas, etiquetándolas como supuestos
académicos. \\
ProFuturo como competidor. & Lo corregí: su web documenta una alianza
con NASCO; lo trato como aliado. \\
Plantilla secciones 9-15 completas. & Las escribí a partir de mis
decisiones: elegí las metas, los buyer persona, el customer journey, los
KPIs y el presupuesto. Devin me ayudó con la estructura, pero el
contenido y las decisiones son míos. \\
\end{longtable}

\hypertarget{prompts-principales-efectivamente-utilizados}{%
\subsubsection{2. Prompts principales efectivamente
utilizados}\label{prompts-principales-efectivamente-utilizados}}

\begin{longtable}[]{@{}
  >{\raggedleft\arraybackslash}p{(\columnwidth - 4\tabcolsep) * \real{0.0045}}
  >{\raggedright\arraybackslash}p{(\columnwidth - 4\tabcolsep) * \real{0.8454}}
  >{\raggedright\arraybackslash}p{(\columnwidth - 4\tabcolsep) * \real{0.1501}}@{}}
\toprule\noalign{}
\begin{minipage}[b]{\linewidth}\raggedleft
Orden
\end{minipage} & \begin{minipage}[b]{\linewidth}\raggedright
Prompt utilizado
\end{minipage} & \begin{minipage}[b]{\linewidth}\raggedright
Contexto y uso
\end{minipage} \\
\midrule\noalign{}
\endhead
\bottomrule\noalign{}
\endlastfoot
\textbf{1} & ``Estoy realizando un proyecto académico en el que debo
seleccionar una organización entre varias alternativas propuestas. Te
voy a proporcionar el documento con las instrucciones del proyecto y las
diferentes opciones disponibles. Actúa como asesor académico y analiza
cuál de las alternativas presenta un mejor encaje con los requisitos del
proyecto. Lee y compara las opciones teniendo en cuenta los objetivos,
requisitos y contenidos solicitados, explicando las ventajas y posibles
dificultades de cada una. Utiliza un lenguaje académico, claro y
didáctico, adecuado para un estudiante universitario. Presenta la
comparación de forma ordenada y finaliza con una síntesis de los
aspectos que debería tener en cuenta para tomar mi decisión. No inventes
información y basa el análisis en los materiales proporcionados.
Responde en español de España.'' & Se aportó el PDF inicial y cuatro
alternativas. La IA comparó su encaje académico y posteriormente yo
elegí NASCO. \\
\textbf{2} & ``He seleccionado NASCO ICT para realizar mi proyecto
académico. Su página web es \url{https://nascoict.org/} y necesito
completar los apartados iniciales del trabajo:''Actividad principal de
la empresa'' y ``Por qué he elegido este proyecto''. Actúa como asesor
académico y analiza la información disponible sobre la organización para
ayudarme a redactar ambos apartados. Utiliza un lenguaje académico,
claro, natural y adecuado para un estudiante universitario. Estructura
la respuesta en los dos apartados indicados, con una extensión
suficiente para explicar las ideas principales sin extenderte
innecesariamente. No inventes información sobre NASCO ICT y, si algún
dato no puede verificarse, indícalo. Responde en español de España.'' &
Se solicitó una presentación inicial de NASCO ICT y una justificación de
la elección del proyecto. \\
\textbf{3} & ``Estoy realizando un dossier académico sobre NASCO ICT,
una ONG sobre la que debo responder varias preguntas relacionadas con su
marca y su entorno. Te proporcionaré los materiales del proyecto y los
contenidos docentes disponibles. El trabajo requiere utilizar
evidencias, realizar un análisis PESTEL y mantener la trazabilidad de
las fuentes. Actúa como investigador y asesor académico. Analiza la
información proporcionada y amplía la investigación necesaria para poder
responder correctamente a las cuestiones del dossier. Utiliza un
lenguaje académico, claro y didáctico, diferenciando los datos objetivos
de las interpretaciones. Organiza la información siguiendo los apartados
y preguntas del dossier e indica las fuentes utilizadas para justificar
las afirmaciones. No inventes datos ni fuentes, prioriza fuentes
oficiales y fiables y respeta el contenido que se haya trabajado en
clase. Responde en español de España.'' & Se amplió la investigación
sobre NASCO ICT teniendo en cuenta las preguntas preliminares, los
contenidos docentes y los requisitos de evidencias, PESTEL y
trazabilidad. \\
\textbf{4} & ``Estoy elaborando un dossier académico sobre NASCO ICT y
necesito investigar su presencia en redes sociales para analizar su
marca y comunicación. Te proporcionaré el PDF completo del proyecto,
aunque en clase únicamente hemos trabajado hasta el Tema 2.1: Marca.
Actúa como investigador especializado en comunicación digital y marca.
Busca y verifica cuáles son las redes sociales oficiales de NASCO ICT y
recopila la información relevante para el análisis de su marca. Utiliza
un lenguaje académico, claro y descriptivo. Presenta los resultados de
forma ordenada, preferiblemente en una tabla indicando la plataforma, la
cuenta oficial y la información relevante. Incluye las fuentes
utilizadas. Prioriza siempre los perfiles oficiales, no confundas
cuentas de terceros con las de la organización y no avances hacia
contenidos correspondientes a temas posteriores al Tema 2.1. No inventes
información y responde en español de España.'' & Se pidió investigar las
redes sociales oficiales de NASCO ICT, limitando el análisis al
contenido de marca trabajado hasta el Tema 2.1. \\
\textbf{5} & ``Necesito elaborar el dossier de mi proyecto académico
sobre NASCO ICT. Hasta este momento, el contenido trabajado en clase
corresponde principalmente a Marca y Entorno, por lo que debo centrar el
trabajo en estos aspectos e incluir un análisis PESTEL. Es importante
utilizar fuentes verificables y mantener la trazabilidad de la
información. Actúa como asesor académico especializado en marketing y
análisis del entorno. Crea el dossier completo en formato Markdown a
partir de los materiales y de la información recopilada, integrando los
apartados de Marca y Entorno y desarrollando el análisis PESTEL. Utiliza
un lenguaje académico, claro, natural y coherente con un trabajo
universitario. Organiza el documento mediante títulos y subtítulos en
Markdown e incluye las fuentes utilizadas. No inventes información ni
fuentes y no desarrolles todavía los apartados de Recursos Humanos,
presupuesto ni competencia, ya que todavía no corresponden al contenido
trabajado en clase. Responde en español de España.'' & Se solicitó crear
el dossier en formato Markdown, incorporando Marca, Entorno y PESTEL y
respetando el límite del temario trabajado. \\
\textbf{6} & ``Ahora debo ampliar el dossier de NASCO ICT con los
apartados de Recursos Humanos, Recursos Económicos y Competencia, ya que
el temario de la asignatura ha avanzado y la entrega final del EA1 los
requiere. Te proporcionaré el dossier actual y los materiales docentes
correspondientes. Actúa como asesor académico especializado en
marketing. Propón una estructura para estos tres apartados a partir de
la información pública disponible sobre NASCO y de las fuentes ya
citadas en el dossier. Utiliza un lenguaje académico, claro y coherente
con el resto del documento. Diferencia los datos objetivos de las
hipótesis y no inventes información sobre la organización. Indica las
fuentes utilizadas y responde en español de España.'' & Se amplió el
dossier con una plantilla de estructura para Recursos Humanos, Recursos
Económicos y Competencia, que yo revisé y completé con mis
decisiones. \\
\textbf{7} & ``Para terminar el EA1 necesito desarrollar los apartados
de Buyer Persona, Customer Journey, Objetivos SMART y KPIs para NASCO
ICT. Te adjunto el dossier actual y los materiales docentes sobre estos
temas. Actúa como asesor académico especializado en inbound marketing.
Propón una plantilla para cada apartado a partir del diagnóstico de
marca y entorno ya realizado, diferenciando hipótesis de datos
verificados. Utiliza un lenguaje académico, claro y coherente con el
resto del documento. No inventes datos internos de NASCO y deja
explícito qué información faltaría confirmar con la organización. Indica
las fuentes utilizadas y responde en español de España.'' & Se generó
una plantilla para Buyer Persona, Customer Journey, Objetivos SMART y
KPIs que yo completé con mis decisiones y mi redacción. \\
\end{longtable}

\begin{center}\rule{0.5\linewidth}{0.5pt}\end{center}

\hypertarget{bibliografuxeda-y-fuentes}{%
\subsection{Bibliografía y fuentes}\label{bibliografuxeda-y-fuentes}}

En páginas sin fecha de actualización verificable utilizo ``s. f.''. La
fecha de consulta no convierte un dato histórico en actual.

\hypertarget{material-docente}{%
\subsubsection{Material docente}\label{material-docente}}

\begin{itemize}
\tightlist
\item
  Iraola Pérez-Tabernero, Marta. \emph{Inbound Marketing: Análisis de
  marca y definición de objetivos}. Unidad 1 de Inbound Marketing e
  Inteligencia Artificial Generativa, Universidad de Deusto. PDF
  aportado: \texttt{Inbound\ Marketing\ e\ IAg\_UD1.pdf}. Referencias
  utilizadas: Tema 2.1, pp.~impresas 39-60; Tema 2.2, pp.~61-103,
  especialmente plantilla PESTEL y ejercicios de pp.~96-99.
\item
  Universidad de Deusto, ALUD. \emph{Del análisis a los objetivos
  medibles de tu estrategia Inbound. Diagnóstico y brújula}. Escenario
  de Aprendizaje 1. PDF de dos páginas aportado por mí.
\item
  Sinek, Simon (2009). \emph{How Great Leaders Inspire Action}. Marco
  conceptual estudiado a través del material docente; no se presenta
  como consulta adicional del vídeo original.
\item
  Kotler, Philip y Keller, Kevin Lane (2016). \emph{Marketing
  Management}, 15.ª edición. Modelo de Producto Ampliado trabajado a
  través del material docente; no afirmo haber consultado el libro
  completo.
\end{itemize}

\hypertarget{marca-y-canales}{%
\subsubsection{Marca y canales}\label{marca-y-canales}}

\begin{itemize}
\tightlist
\item
  \textbf{\href{https://nascoict.org/}{F01}} NASCO Feeding Minds (s.
  f.). \emph{Portada en español}. Identidad, mensaje principal, cifra de
  aulas y enlaces a redes.
\item
  \textbf{\href{https://nascoict.org/quienes-somos/}{F02}} NASCO Feeding
  Minds (s. f.). \emph{Quiénes somos}. Misión, historia y cifras
  publicadas.
\item
  \textbf{\href{https://nascoict.org/wp-content/uploads/MEMORIA-2024.pdf}{F03}}
  NASCO Feeding Minds. \emph{Memoria anual 2024}. Documento de actividad
  correspondiente a 2024; distingo sus datos de la situación actual.
  Páginas clave: 3-5, 9-10 y 13-15.
\item
  \textbf{\href{https://nascoict.org/aulas-informaticas/}{F04}} NASCO
  Feeding Minds (s. f.). \emph{Aulas informáticas}. Oferta educativa,
  cifras y colaboración tecnológica.
\item
  \textbf{\href{https://nascoict.org/hazte-socio/}{F05}} NASCO Feeding
  Minds (s. f.). \emph{Hazte socio}. Modalidades, importes, formulario y
  condiciones anunciadas.
\item
  \textbf{\href{https://nascoict.org/transparencia/}{F06}} NASCO Feeding
  Minds (s. f.). \emph{Transparencia}. Índice documental 2017-2024.
\item
  \textbf{\href{https://nascoict.org/blog/}{F07}} NASCO Feeding Minds
  (s. f.). \emph{Blog}. Primera página recuperada; entradas visibles
  hasta el 31/05/2023.
\item
  \textbf{\href{https://www.instagram.com/nasco_ict/}{F08}} NASCO
  Feeding Minds (s. f.). \emph{Perfil de Instagram @nasco\_ict}.
  Contadores y biografía recuperados de metadatos públicos.
\item
  \textbf{\href{https://www.facebook.com/NascoICT/}{F09}} NASCO Feeding
  Minds (s. f.). \emph{Página de Facebook NascoICT}. Contador y
  descripción recuperados de metadatos públicos.
\item
  \textbf{\href{https://www.linkedin.com/company/nasco-feeding-minds/}{F10}}
  NASCO Feeding Minds (s. f.). \emph{Perfil corporativo de LinkedIn}.
  Identificación del canal; las cifras y publicaciones se contrastaron
  mediante F12-F14.
\item
  \textbf{\href{https://twitter.com/nasco_ict/}{F11}} NASCO Feeding
  Minds (s. f.). \emph{Perfil de X @nasco\_ict}. Contadores y muestra
  pública recuperada.
\item
  \textbf{\href{https://www.linkedin.com/posts/nasco-feeding-minds_nascofoundation-educaci\%C3\%B3ndigital-impactosocial-activity-7500847528586846208-5q6E}{F12}}
  NASCO Feeding Minds (02/09/2026). \emph{La transformación digital
  comienza con una educación básica sólida}. Publicación de LinkedIn.
\item
  \textbf{\href{https://www.linkedin.com/posts/nasco-feeding-minds_nascofoundation-feedingminds-digitaleducation-activity-7465326634426843136-PKBW}{F13}}
  NASCO Feeding Minds (27/05/2026). \emph{Inauguración del aula de
  educación digital en Tuna}. Publicación de LinkedIn.
\item
  \textbf{\href{https://www.linkedin.com/posts/nasco-feeding-minds_nasco-x-deloitte-activity-7464593630444306432-8i9r}{F14}}
  NASCO Feeding Minds (25/05/2026). \emph{Un gran paso hacia el futuro:
  NASCO y Deloitte}. Publicación de LinkedIn.
\item
  \textbf{\href{https://nascoict.org/aviso-legal-y-politica-de-privacidad/}{F29}}
  NASCO Feeding Minds (s. f.). \emph{Aviso legal y política de
  privacidad}. Fuente contextual sobre la información que comunica la
  organización; no constituye una auditoría de cumplimiento.
\item
  \textbf{\href{https://nascoict.org/en/}{F30}} NASCO Feeding Minds (s.
  f.). \emph{Portada en inglés}. Fuente de la afirmación sobre
  subvenciones gubernamentales pendiente de aclaración frente al alcance
  de F03.
\end{itemize}

\hypertarget{entorno}{%
\subsubsection{Entorno}\label{entorno}}

\begin{itemize}
\tightlist
\item
  \textbf{\href{https://knowledge4policy.ec.europa.eu/publication/oecd-preliminary-official-development-assistance-levels-2025_en}{F15}}
  Comisión Europea, Knowledge4Policy (2026). \emph{OECD Preliminary
  official development assistance levels in 2025}. Síntesis
  institucional de datos preliminares de la OCDE; dato de 2025.
\item
  \textbf{\href{https://www.boe.es/buscar/act.php?id=BOE-A-2023-4512}{F16}}
  BOE (21/02/2023). \emph{Ley 1/2023, de 20 de febrero, de Cooperación
  para el Desarrollo Sostenible y la Solidaridad Global}. Texto
  consolidado; artículo 1.6.
\item
  \textbf{\href{https://www.ine.es/dyngs/Prensa/es/ECV2025.htm}{F17}}
  INE (05/02/2026). \emph{Encuesta de Condiciones de Vida. Año 2025.
  Resultados definitivos}.
\item
  \textbf{\href{https://www.worldbank.org/en/news/press-release/2026/08/26/ghana-must-turn-strong-recovery-into-quality-jobs-and-lasting-growth}{F18}}
  Banco Mundial (26/08/2026). \emph{Ghana Must Turn Strong Recovery into
  Quality Jobs and Lasting Growth}. Comunicado sobre la décima
  actualización económica de Ghana.
\item
  \textbf{\href{https://www.aefundraising.org/perfil-donante-2024-que-impulsa-a-los-espanoles-a-colaborar/}{F19}}
  Asociación Española de Fundraising (04/12/2024). \emph{Perfil donante
  2024: ¿Qué impulsa a los españoles a colaborar?} Resumen sectorial del
  estudio realizado con Tuinkel; no he revisado su microdato.
\item
  \textbf{\href{https://www.ine.es/dyngs/Prensa/es/TICH2025.htm}{F20}}
  INE (20/11/2025). \emph{Encuesta sobre Equipamiento y Uso de
  Tecnologías de la Información y Comunicación en los Hogares. Año
  2025}.
\item
  \textbf{\href{https://www.itu.int/en/ITU-D/Environment/Pages/Publications/The-Global-E-waste-Monitor-2024.aspx}{F21}}
  ITU/UNITAR (20/03/2024). \emph{The Global E-waste Monitor 2024}.
  Síntesis oficial de resultados; cifras de residuos correspondientes a
  2022.
\item
  \textbf{\href{https://www.worldbank.org/en/news/press-release/2022/11/01/ghana-can-turn-climate-challenges-into-opportunities-for-resilient-and-sustainable-growth-says-new-world-bank-group-report}{F22}}
  Banco Mundial (01/11/2022). \emph{Ghana Can Turn Climate Challenges
  into Opportunities for Resilient and Sustainable Growth}. Comunicado
  del diagnóstico climático; lo utilizo como contexto estructural, no
  como recuento de daños de 2026.
\item
  \textbf{\href{https://www.aepd.es/prensa-y-comunicacion/notas-de-prensa/aepd-actualiza-guia-cookies-para-adaptarla-a-nuevas-directrices-cepd}{F23}}
  AEPD (11/07/2023). \emph{La AEPD actualiza su Guía sobre el uso de
  cookies para adaptarla a las nuevas directrices del Comité Europeo de
  Protección de Datos}.
\item
  \textbf{\href{https://sede.agenciatributaria.gob.es/Sede/ayuda/manuales-videos-folletos/manuales-practicos/irpf-2024/c16-deducciones-generales-cuota/deducciones-donativos-otras-aportaciones/donativos-realizados-entidades-incluidas/general-donativos-donaciones-entidades-beneficiarias-mecenazgo/porcentajes-deduccion.html}{F24}}
  AEAT (actualización 10/04/2025). \emph{Manual práctico de Renta 2024:
  porcentajes de deducción por donativos}. Referencia del ejercicio
  2024, no garantía automática para otros ejercicios o territorios.
\item
  \textbf{\href{https://trends.zorraquino.com/informes/consumer-trends/2026/}{F25}}
  Zorraquino. \emph{Consumer Trends 2026}. Presentación pública de la
  edición; fecha exacta de actualización de la página no verificada.
\item
  \textbf{\href{https://trends.zorraquino.com/informes/consumer-trends/2026/brain-rot.html}{F26}}
  Zorraquino. \emph{Consumer Trends 2026: Brain rot}. Desarrollo público
  de la tendencia 6.
\item
  \textbf{\href{https://trends.zorraquino.com/informes/consumer-trends/2026/cognitiveco-agents.html}{F27}}
  Zorraquino. \emph{Consumer Trends 2026: Cognitive co-agents}.
  Desarrollo público de la tendencia 9.
\item
  \textbf{\href{https://trends.zorraquino.com/informes/consumer-trends/2026/everythingcampaignable.html}{F28}}
  Zorraquino. \emph{Consumer Trends 2026: Everything campaignable}.
  Desarrollo público de la tendencia 4.
\end{itemize}

\hypertarget{alternativas-y-colaboraciones-observadas}{%
\subsubsection{Alternativas y colaboraciones
observadas}\label{alternativas-y-colaboraciones-observadas}}

\begin{itemize}
\tightlist
\item
  \textbf{\href{https://teso.org.es/informacion/donar-material/}{F31}}
  TeSo. \emph{Donaciones de material}. Requisitos, cuidado de datos y
  procedimiento de aportación de equipos; consulta 10/10/2026.
\item
  \textbf{\href{https://teso.org.es/entrega-del-proyecto-oyera-d-a-school-de-ghana/}{F32}}
  TeSo (30/05/2023). \emph{Entrega del proyecto Oyera D/A School de
  Ghana}. Ejemplo concreto de dotación tecnológica; consulta 10/10/2026.
\item
  \textbf{\href{https://www.entreculturas.org/hazte-socio/}{F33}}
  Entreculturas. \emph{Hazte socio/a}. Página de colaboración de
  particulares; consulta 10/10/2026.
\item
  \textbf{\href{https://www.entreculturas.org/hazte-socio-empresas/}{F34}}
  Entreculturas. \emph{Hazte socio como empresa}. Página de colaboración
  empresarial; consulta 10/10/2026.
\item
  \textbf{\href{https://www.worldreader.org/ways-to-donate/}{F35}}
  Worldreader. \emph{Ways to Donate}. Opciones públicas de aportación;
  consulta 10/10/2026.
\item
  \textbf{\href{https://www.worldreader.org/now/worldreader-and-compassion-international-join-forces/}{F36}}
  Worldreader (02/07/2022). \emph{Partnership Announcement: Worldreader
  and Compassion International Ghana}. Programa anunciado en Ghana;
  consulta 10/10/2026.
\item
  \textbf{\href{https://profuturo.education/noticias/profuturo-nasco-feeding-minds-educacion-digital-ghana/}{F37}}
  ProFuturo. \emph{ProFuturo y NASCO Feeding Minds: unidos por la
  educación digital en Ghana}. Fuente de su relación de colaboración;
  consulta 10/10/2026.
\end{itemize}

\end{document}
